 \documentclass[11pt,twoside]{amsart}
\usepackage[dvipsnames]{xcolor}
\usepackage{hyperref}
\hypersetup{colorlinks=true, linkcolor=Blue, citecolor=Maroon}
\usepackage{amsthm, amsmath, amscd, amssymb,centernot,txfonts}
\usepackage{tikz-cd}
\usepackage{lipsum}
\usepackage{cancel} 
\usepackage{multicol}
\usepackage{amsfonts}
\usepackage{mathrsfs}
\usepackage[all]{xy}
\usepackage[left=2.5cm,top=2.5cm,bottom=3cm,right=2.5cm]{geometry}
\usepackage{dirtytalk}
\usepackage{mathtools}
\usetikzlibrary{shapes,arrows}
\usepackage{etoolbox}
\usepackage{dynkin-diagrams}
\usepackage{enumitem}
\usepackage{chngpage}
\usepackage{adjustbox}
\usepackage[normalem]{ulem}
\usepackage[utf8]{inputenc}
\usepackage{fourier} 
\usepackage{array}
\usepackage{makecell}
\usepackage{comment}
%\usepackage[cmtip,all]{xy}
\usepackage{cancel}
\usepackage{cleveref}
\newcommand{\longsquiggly}{\xymatrix{{}\ar@{~>}[r]&{}}}




\setlength{\headheight}{15.2pt}

\renewcommand\theadalign{bc}
\renewcommand\theadfont{\bfseries}
\renewcommand\theadgape{\Gape[4pt]}
\renewcommand\cellgape{\Gape[4pt]}



\setcounter{tocdepth}{1}

\setlength\parskip{0in}
\setlength\parindent{0.2in}


\theoremstyle{plain}
\numberwithin{equation}{section}
\newtheorem{theorem}{Theorem}[section]
\newtheorem{proposition}[theorem]{Proposition}
\newtheorem{lemma}[theorem]{Lemma}
\newtheorem{corollary}[theorem]{Corollary}
\newtheorem{conjecture}[theorem]{Conjecture}
\newtheorem{claim}[theorem]{Claim}




\newtheorem{observation}[theorem]{Observation}


\theoremstyle{definition}
\newtheorem{remark}[theorem]{Remark}
\newtheorem{notation}[theorem]{Notation}
\newtheorem{example}[theorem]{Example}
\newtheorem{question}[theorem]{Question}
\newtheorem{set-up}[theorem]{Set-up}
\newtheorem{definition}[theorem]{Definition}

\newcommand{\te}{\otimes}

\newcommand{\myeq}{\overset{\mathrm{duality}}{=\joinrel=}}
\newcommand*{\QEDA}{\hfill\ensuremath{\blacksquare}}
\newcommand*{\QEDB}{\hfill\ensuremath{\square}}

\tikzstyle{decision} = [diamond, draw, , 
    text width=4.5em, text badly centered, node distance=3cm, inner sep=0pt]
\tikzstyle{block} = [rectangle, draw, , 
    text width=10em, text centered, rounded corners, minimum height=2em]
\tikzstyle{block1} = [rectangle, draw, , 
    text width=5em, text centered, rounded corners, minimum height=2em]    
\tikzstyle{line} = [draw, -latex']
\tikzstyle{cloud} = [draw, ellipse,, node distance=3cm,
    minimum height=2em]
\begin{document}

\title[Extendability of $K3$ surfaces without Gaussian maps and classification of Fano threefolds]{Extendability of $K3$ surfaces without Gaussian maps and classification of Fano threefolds}


%\author[P. Bangere]{Purnaprajna Bangere}
%\address{Department of Mathematics, University of Kansas, Lawrence, USA}
%\email{purna@ku.edu}

%\author[F.J. Gallego]{Francisco Javier Gallego}
%\address{Departamento de \'Algebra, Geometr\'ia y Topolog\'ia and Instituto de Matem\'atica Interdisciplinar,
%Universidad Complutense de Madrid, Spain}
%\email{gallego@mat.ucm.es}

%\author[J. Mukherjee]{Jayan Mukherjee}
%\address{Department of Mathematics, Oklahoma State University, Stillwater, USA}
%\email{jayan.mukherjee@okstate.edu}

\subjclass[2020]{14B05, 14B10, 14D06, 14D15, 14D20, 14E30, 14J10, 14J45}
\keywords{compactification of moduli spaces, simple normal crossings, semi-log-canonical singularities, smoothings, deformations of morphisms, multiple structures, Fano varieties, Calabi-Yau varieties, varieties of general type, moduli of varieties of general type.}

\maketitle
\begin{abstract}

%In this article we show the existence of simple normal crossing varieties $V = Y_1 \bigcup_E Y_2 \subset \mathbb{P}^N$, where each $Y_i \hookrightarrow \mathbb{P}^N$ is a Hirzebruch surface embedded by the complete linear series of a very ample line bundle $aC_0+bf$ that intersect along an anticanonical elliptic curve $E$. We show further that these snc varieties are smoothable in $\mathbb{P}^N$ into smooth $K3$ surfaces. For a fixed $a$ and $b$, these polarized type II degenerations form a smooth locus of codimension $6$ inside the corresponding Hilbert scheme. We show that when $Y_i = \mathbb{F}_0$ and $a = 1$, a general such $V$, which is a union of scrolls, has well-defined and semistable syzygy points $\textrm{Syz}_{0,q}(V)$ for $q \geq 2$ and $\textrm{Syz}_{p,2}(V)$ for $p \leq \lfloor \displaystyle\frac{g-1}{2} \rfloor$ where $g$ is the sectional genus of the smoothing $K3$ fiber. We draw similar conclusions for the syzygy points of the general canonical binary curve sections of $V$.
 

\end{abstract}

\section{Introduction}


\section{General theorems on projective smoothing of simple normal crossings }\label{main results}

\begin{comment}
We assume throughout that we work over an algebraically closed field of characteristic zero. 

\begin{lemma}\label{Hilbert polynomial}
Let $Y$ be a smooth projective scheme embedded (possibly degenerately) inside a projective space 
 $\mathbb{P}^N$ and let $D \in \mid L^{-1} \mid$ be a smooth divisor in $Y$. Suppose there exist
\begin{itemize}
    \item[(1)] A scheme $V$ inside $\mathbb{P}^N$ which is the (scheme-theoretic) union $Y_1 \bigcup Y_2$ of two smooth schemes inside $\mathbb{P}^N$ such that both $Y_1$ and $Y_2$  belong to the same Hilbert scheme as $Y$ and $Y_1 \bigcap Y_2 = D$ (scheme-theoretic intersection) as a subscheme of $\mathbb{P}^N$. Then,  
    $$ \chi(\mathscr{O}_V(t)) = \chi(\mathcal{O}_{Y_1}(t)) +  \chi(\mathcal{O}_{Y_2}(t)) - \chi(\mathcal{O}_{D}(t)) $$
      
     and, for $t>>0$,
     $$ \chi(\mathscr{O}_V(t)) = 2\chi(\mathcal{O}_{Y}(t))-\chi(\mathcal{O}_{D}(t)).$$
    \item[(2)] A ribbon $\widetilde{Y}$ on $Y$ with conormal bundle $L$ embedded inside $\mathbb{P}^N$. Then, for $t>>0$, $$ \chi(\mathscr{O}_{\widetilde{Y}}(t)) = 2\chi(\mathcal{O}_{Y}(t))-\chi(\mathcal{O}_{D}(t)).$$
\end{itemize}
In particular, the Hilbert polynomials of $V$ and $\widetilde{Y}$ are the same. 
\end{lemma}

\begin{remark}\label{ideal of a ribbon}
(see the proof of \cite[Lemma 1.4]{GP97}; see also \cite{HV})
Let $Y \hookrightarrow \mathbb{P}^N$ be an embedding of a smooth variety inside a projective space. Let $\mathcal{I}$ denote the ideal sheaf of $Y$. Let $\widetilde{Y}$ denote an embedded ribbon on $Y$ with conormal bundle $L$ corresponding to a surjective homomorphism $\Bar{\psi} \in \textrm{Hom}(\mathcal{I}/\mathcal{I}^2, L) = H^0(N_{Y/\mathbb{P}^N} \otimes L)$. Then the ideal sheaf $\mathcal{I}_{\widetilde{Y}}$ of $\widetilde{Y}$ inside $\mathbb{P}^N$ is given by $\textrm{Ker}\psi$ where $\psi$ is the composition of $\psi: \mathcal{I} \to \mathcal{I}/\mathcal{I}^2 \to L$
\end{remark}


\begin{lemma}\label{deformations fixing a subscheme}
Let $D \hookrightarrow Y \hookrightarrow \mathbb{P}^N$ be the embedding of smooth subvariety $D$ and variety  $Y$ inside a projective space over an algebraically closed field $\mathbf k$ of characteristic zero. Let $\textrm{Art}_\mathbf k$ denote the category of artinian $\mathbf k$-algebras. Let $F: \textrm{Art}_\textbf k \to \textrm{Sets}$ be the functor as follows: For an Artinian $\mathbf k-$algebra $A$,
$$F_{D,Y}(A) = \left\{ D \times \textrm{Spec}(A) \hookrightarrow Y_A \hookrightarrow \mathbb{P}^N \times \textrm{Spec}(A) \right\}$$
where $Y_A \hookrightarrow \mathbb{P}^N \times \textrm{Spec}(A)$ is an embedded deformation of $Y$ inside $\mathbb{P}^N$ over $\textrm{Spec}(A)$. Then 
\begin{itemize}
    \item[(1)] $F_{D,Y}(\mathbf k[\epsilon]/\epsilon^2) = H^0(N_{Y/\mathbb{P}^N} \otimes \mathcal{I}_{D/Y})$ where $\mathcal{I}_{D/Y}$ is the ideal sheaf of $D$ inside $Y$.
    \item[(2)] $H^1(N_{Y/\mathbb{P}^N} \otimes \mathcal{I}_{D/Y})$ is an obstruction space for the functor $F_{D,Y}$. 
\end{itemize}
\end{lemma}

\noindent\textit{Proof.} The functor $F_{D,Y}$ parametrizing subschemes $Y$ of $\mathbb P^N$ that contain $D$, is a natural generalization of the Hilbert scheme (see \cite[p. 230]{Ser}). Then the result follows as a natural generalization of \cite[Theorem VI-29]{EH} or the analogous result in \cite[Proposition 3.2.6]{Ser} or \cite{H10}. \QEDA

\begin{remark}\label{flag hilbert schemes}
Let $i: Y \hookrightarrow \mathbb{P}^N$ be an embedding of a smooth variety inside a projective space. Let $D \in |L^{-1}|$ be a smooth divisor. Then there is an induced embedding of $j: D \hookrightarrow \mathbb{P}^N$. Let $p_1$ and $p_2$ denote the Hilbert polynomials of $Y \hookrightarrow \mathbb{P}^N$ and $D \hookrightarrow \mathbb{P}^N$ respectively. Let $\mathcal{H}(p_1, p_2)$ denote the flag Hilbert scheme representing the flag-Hilbert functor as defined in \cite{Ser}, Section $4.5.1$ or \cite{KL81}. Closed points of $\mathcal{H}(p_1, p_2)$ parameterize pairs $(D' \hookrightarrow Y' \hookrightarrow \mathbb{P}^N)$ where the subschemes $Y' \hookrightarrow \mathbb{P}^N$ and $D' \hookrightarrow \mathbb{P}^N$ have Hilbert polynomials $p_1$ and $p_2$ respectively. Let $\mathcal{D}$ denote the Hilbert scheme parameterizing subschemes $(D' \hookrightarrow \mathbb{P}^N)$ with Hilbert polynomial $p_2$. There is a projection map $p: \mathcal{H}(p_1, p_2) \to \mathcal{D}$ whose fibre $\mathcal{F}_D$ at a subscheme $(D \hookrightarrow \mathbb{P}^N)$ consists of all those subschemes $Y' \hookrightarrow \mathbb{P}^N$ with Hilbert polynomial $p_1$ such that $D \subset Y'$. 
%Consider the functor $G$ represented by the fibre of $p: \mathcal{H}(p_1,p_2) \to \mathcal{D}$ at the point $(D \hookrightarrow \mathbb{P}^N) \in \mathcal{D}(k)$. 
%Now consider the restriction of the functor $G$ at the point $(D \hookrightarrow Y \hookrightarrow \mathbb{P}^N) \in \mathcal{H}(p_1,p_2)(k)$ of the fibre to the category of local Artinian $k$- algebras. 
It follows from the definition of flag-Hilbert functor $\mathcal{H}(p_1,p_2)$ given in \cite{Ser} or \cite{KL81}, that $F_{D,Y}$ defined in Lemma \ref{deformations fixing a subscheme} is represented by the pointed scheme consisting of  $(\mathcal{F}_D, (D \hookrightarrow Y \hookrightarrow \mathbb{P}^N))$. In particular the dimension of the fibre at $(D \hookrightarrow Y \hookrightarrow \mathbb{P}^N)$ is $H^0(N_{Y/\mathbb{P}^N} \otimes L)$ and the fibre is smooth at the point $(D \hookrightarrow Y \hookrightarrow \mathbb{P}^N)$ if $H^1(N_{Y/\mathbb{P}^N} \otimes L) = 0$. Let $\mathcal F_{D,Y}$ denote the unique irreducible component of $\mathcal F_D$ containing the point $(D \hookrightarrow Y\hookrightarrow \mathbb{P}^N)$. Closed points parameterized by $\mathcal{F}_{D,Y}$ correspond to deformations over an irreducible curve of the subvariety $Y$ inside $\mathbb{P}^N$ which keeps $D$ fixed.     
\end{remark}


\color{black}


















\begin{definition}
    Let $Y_1$ and $Y_2$ be two smooth subvarieties of $\mathbb P^N$ and let $D$ be a smooth divisor on both $Y_1$ and $Y_2$. If $D$ is the scheme-theoretical intersection of we denote the union $Y_1 \bigcup Y_2$ as $Y_1 \bigcup_D Y_2$.
\end{definition}

 \textcolor{red}{we need to refer our paper [BGM24b], Theorem 2.7 here}

\color{black}
\begin{theorem}\label{main}
Let $Y$ be a smooth projective variety embedded in $\mathbb{P}^N$. Let $L$ be a line bundle on $Y$ and  
$D \in \mid L^{-1} \mid$  be a smooth divisor. Suppose there exists a nowhere vanishing section $\zeta$ in $H^0(N_{Y/\mathbb{P}^N} \otimes L)$ which corresponds to 
a ribbon 
$\widetilde{Y}$, embedded in $\mathbb{P}^N$, supported on $Y$. Assume 
that $Y$ is an unobstructed subscheme of $\mathbb{P}^N$, and that the functor $F_{D,Y}$ is unobstructed or equivalently the fibre of the map $p: \mathcal{H}(p_1,p_2) \to \mathcal{D}$ over the point $(D \hookrightarrow \mathbb{P}^N)$ is smooth at the point $(D \hookrightarrow Y \hookrightarrow \mathbb{P}^N)$. Then there exists a flat family of subschemes $\mathcal{Y} \to T$ of $\mathbb{P}^N_T$  over a smooth irreducible algebraic curve $T$ such that 
\begin{itemize}
    \item[(1)]  $\mathcal{Y}_0$ is the subscheme $\widetilde{Y}$. 
    \item[(2)] $\mathcal{Y}_t = Y_{1t} \bigcup Y_{2t}$, where $Y_{1t}$ is the subscheme $Y$, $Y_{2t}$ is a non trivial deformation in $\mathbb{P}^N_T$ of 
    the subscheme $Y$ that fixes the subscheme $D$
    (i.e., 
    for any $t \in T, t \neq 0$, $Y_{2t}$ and $Y$ are different subschemes of $\mathbb{P}^N_T$)
    and $Y_{1t} \bigcap Y_{2t} \cong D$.
\end{itemize}
\end{theorem}


\begin{remark}\label{blow-up of a ribbon}
    Let $\widetilde{Y}$ be a ribbon over $Y$ with conormal bundle $L$. Let $D \in |L^{-1}|$ denote a smooth divisor on $Y$. Then $D$ defines a Weil-divisor on $\widetilde{Y}$. Choosing a section $s$ corresponding to $D$, we have a map $H^1(T_Y \otimes L) \to H^1(T_Y)$ obtained by multiplication by the section $s$. Note that $\widetilde{Y}$ represents an element of $H^1(T_Y \otimes L)$ and under the above map the image of $\widetilde{Y}$ is the blow-up of $\widetilde{Y}$ along the Weil divisor $D$, which is a ribbon with conormal bundle $L \otimes D = \mathcal{O}_Y$ (see \cite[Theorem 1.9]{BE95}) . A ribbon with conormal bundle $\mathcal{O}_Y$ is by the definition of a ribbon (see Definition \ref{def of a ribbon}) a first order deformation of $Y$, so indeed the image lies $H^1(T_Y)$. Now in the embedded setting, we have a commutative diagram
    \[
    \begin{tikzcd}
        H^1(N_{Y/\mathbb{P}^N} \otimes L) \arrow[d] \arrow[r, "\lambda"] & H^1(N_{Y/\mathbb{P}^N}) \arrow[d] \\
        H^1(T_Y \otimes L) \arrow[r] & H^1(T_Y) \\
    \end{tikzcd}
    \]
where the vertical maps are the forgetful maps. So for the embedded ribbon $\widetilde{Y} \hookrightarrow \mathbb{P}^N$, the image $\lambda(\widetilde{Y})$ is abstractly the blow-up of $\widetilde{Y}$ along $D$ and hence is itself an embedded (in $\mathbb{P}_{\Delta}^N$) first order deformation of $Y$.
\end{remark}



\begin{definition}
    Let $Y$ be a smooth projective variety. A smooth projective variety $Y'$ is deformation equivalent to $Y$ if there exists an irreducible curve $T$ and a flat family $\mathcal{Y}$ such that $\mathcal{Y}_t = Y$ and $\mathcal{Y}_{t'} = Y'$ for two different points $t$ and $t'$. 
\end{definition}

In the rest of the article we will follow the conventions and notations set up in Lemma \ref{deformations fixing a subscheme} and Remark \ref{flag hilbert schemes}. Let the fibre of the map $p: \mathcal{H}(p_1,p_2) \to \mathcal{D}$ over the point $(D \hookrightarrow \mathbb{P}^N)$ be denoted by $\mathcal{F}_D$. Let  $\mathcal F_{D,Y}$ be the union of irreducible components of $\mathcal F_D$ containing the point $(D \hookrightarrow Y\hookrightarrow \mathbb{P}^N)$. Recall from Lemma \ref{deformations fixing a subscheme} that $F_{D,Y}$ is functor of embedded deformations of $Y$ inside $\mathbb{P}^N$ keeping $D$ fixed. Hence if $F_{D,Y}$ is unobstructed, the point $(D \hookrightarrow Y\hookrightarrow \mathbb{P}^N)$ is a smooth point of the fibre $\mathcal{F}_D$ and in particular, $\mathcal{F}_{D,Y}$ is irreducible. 





\textcolor{red}{We need to refer to [BGM24b], Theorem 2.10 here.}


\begin{theorem}\label{intrinsic}
    Let $Y$ be a smooth projective variety and $D \in |L^{-1}|$ be a smooth divisor. Assume there exists an embedding $i: Y \hookrightarrow \mathbb{P}^N$ such {that the} following conditions hold{:}
    \begin{itemize}
        \item[(a)] There exists a  {nowhere vanishing section} {$\nu$} in $H^0(N_{Y/\mathbb{P}^N} \otimes L)$. 
        \item[(b)] The variety {$Y$ is unobstructed in $\mathbb P^N$.}
        \item[(c)] The functor $F_{D,Y}$ is unobstructed  
         (note that this holds if $H^1(N_{Y/\mathbb{P}^N} \otimes L) = 0$). 
        \item[(d)] {The embedded ribbon}  $\widetilde{Y}$ {corresponding to $\nu$} is smoothable in $\mathbb{P}^N$.
        \item[(e)] The Hilbert point of $\widetilde{Y}$ in $\mathbb{P}^N$ is smooth. 
    \end{itemize}
Then:
\begin{itemize}
    \item[(1)] 
   If $Y_1$ is the subvariety $Y$ of $\mathbb P^N$ and $Y_2$ is a subvariety of $\mathbb P^N$ that corresponds to a general element of 
    $\mathcal{F}_{D,Y}$, then {the scheme theoretic intersection of $Y_1 and Y_2$ is $D$ and} 
{$V = Y_1 \bigcup_D Y_2$} is  smoothable inside $\mathbb{P}^N$. The Hilbert point of $V$ inside $\mathbb{P}^N$ is smooth. Further the closure inside the Hilbert scheme of the locus of varieties of the form $V = Y_1 \bigcup_D Y_2$ contains the ribbon $\widetilde{Y}$. 
\end{itemize}

{Let $N_{V/\mathbb{P}^N}'$ denote the equisingular 
normal sheaf of $V$ inside $\mathbb{P}^N$ and assume that, in addition to conditions (a), (b), (c), (d) and (e),  the cohomology group $H^1(N_{V/\mathbb{P}^N}')$ vanishes. Then:} 
\begin{itemize}
    \item[(2)] The locally trivial deformations of $V$ inside $\mathbb{P}^N$ are unobstructed
    and form a subspace in the tangent space of the Hilbert scheme at $[V]$ of codimension $h^0(L^{-2}|_D)$. 
    Furthermore, this subspace is the tangent space of an  irreducible locus of the Hilbert scheme, which has codimension $h^0(L^{-2}|_D)$ and is smooth at $[V]$. The singularities of the subschemes parameterized by this locus are  normal crossing singularities;  {in particular, they are non-normal and semi log canonical. They are analytically isomorphic to $(x_1^2+x_2^2 = 0 \subset \mathbb{C}^N)$.}
\end{itemize}



\end{theorem}



\begin{theorem}\label{theorem.irreducible}
Let $Y$ and $D$ satisfy all the hypotheses of Theorem~\ref{intrinsic}. Assume furthermore that hypotheses (a) and (d) hold for a general element of  $\mathcal{F}_{D,Y}$ ({note this holds if $H^1(N_{Y/\mathbb{P}^N} \otimes L) = 0$, see Remark~\ref{remark.generality}}). {Then for a general pair $(Y_1',Y_2')$ in  $\mathcal{F}_{D,Y} \times \mathcal{F}_{D,Y}$, the union of $Y_1'$ and $Y_2'$ is a normal crossing subvariety {$V=Y_1' \bigcup_D Y_2'$} of $\mathbb{P}^N$, which}
    is smoothable. Such subvarieties $V$ form  an irreducible, flat family. In particular, if $\mathcal{F}_{D}$ is irreducible  {and hypotheses (a) and (d) are fulfilled by a general element of  $\mathcal{F}_{D}$, then the above holds for a general pair $(Y_1',Y_2')$ in  $\mathcal{F}_{D} \times \mathcal{F}_{D}$}.
\end{theorem}
    



\begin{remark}\label{remark.generality}
Let $\mathcal{F}_{D}$ be irreducible
\begin{enumerate}
    \item A sufficient condition for hypotheses (a) and (d) to hold for a general element of $\mathcal{F}_{D}$ is $h^0(N_{\mathcal Y_t/\mathbb P^N} \otimes  L)$ to be constant for general elements of $\mathcal{F}_{D}$. 
    \item A sufficient condition for $h^0(N_{\mathcal Y_t/\mathbb P^N} \otimes  L)$ to be constant for general elements of $\mathcal{F}_{D}$ is $h^1(N_{\mathcal Y_t/\mathbb P^N} \otimes  L)$ to be $0$. 
\end{enumerate}
\end{remark}



\section{Projective degenerations of $K3$ surfaces into a union of Hirzebruch surfaces}

\textcolor{red}{We will now prove a degeneration result that shows that K3 surfaces in projective space degenerates into union of Hirzebruch surfaces embedded by an arbitrary very ample linear series. This result recovers and generalizes the result of [CLM] where they showed K3 surfaces degenerates into union of $Y = \mathbb{F}_0$ embedded as scrolls.}

\begin{theorem}\label{snc to carpet} Let $Y = \mathbb{F}_e \hookrightarrow \mathbb{P}^N$ be an embedding of the Hirzebruch surface $\mathbb{F}_e$, with $e \leq 2$, induced by the complete linear series of the very ample line bundle $|aC_0+bf|$ (hence $b \geq ae+1$). Let $i: \mathbb{P}^N \to \mathbb{P}^M$ be a linear embedding with $M \geq N$ and consider the composed (possibly degenerate) embedding $Y \hookrightarrow \mathbb{P}^M$. Let $E \in |\omega_Y^{-1}|$ denote a smooth elliptic curve. Then there exists a flat family of subschemes $\mathcal{Y} \to T$ of $\mathbb{P}^M_T$ over a smooth irreducible algebraic curve $T$ such that 
\begin{itemize}
    \item[(1)]  $\mathcal{Y}_0$ is an embedded $K3$ carpet $\widetilde{Y} \hookrightarrow \mathbb{P}^M$ extending the embedding $Y \hookrightarrow \mathbb{P}^M$.  
    \item[(2)] $\mathcal{Y}_t = Y_{1t} \bigcup Y_{2t}$, where $Y_{1t}$ is the subvariety $Y$, $Y_{2t}$ is a non trivial deformation in $\mathbb{P}^M_T$ of the subvariety $Y$ that fixes the subvariety $E$ (i.e., for any $t \in T, t \neq 0$, $Y_{2t}$ and $Y$ are different subvarieties of $\mathbb{P}^M_T$)
    and $Y_{1t} \bigcap Y_{2t} \cong E$. 
\end{itemize}
\end{theorem}

\noindent\textit{Proof.} $K3$ carpets $\widetilde{Y} \hookrightarrow \mathbb{P}^M$ that extends the embedding $Y \hookrightarrow \mathbb{P}^M$ correspond to nowhere vanishing sections $\zeta \in H^0(N_{Y/\mathbb{P}^M} \otimes K_Y)$. If $a = 1$, then the existence of such sections follow from \cite[Theorem $1.7$]{GP97}. If $a \geq 2$, then the existence follows from Theorem \cite[Theorem $3.1$]{PMR21}. We show that $Y$ is unobstructed in $\mathbb{P}^M$. Consider the Euler sequence of $\mathbb{P}^M$ pulled back to $Y$, $$0 \to \mathcal{O}_Y \to \mathcal{O}_Y(aC_0+bf)^{M+1} \to T_{\mathbb{P}^M}|_Y \to 0$$
Note that $H^2(\mathcal{O}_Y) = 0$ and $H^1(\mathcal{O}_Y(aC_0+bf)) = 0$ (since $b \geq ae+1$). This implies that $H^1(T_{\mathbb{P}^M}|_Y) = 0$. Combined with the fact that $H^2(T_Y) = 0$ (by \cite[Section $(2)$]{Sei92}), we have that $H^1(N_{Y/\mathbb{P}^M}) = 0$ and hence $Y$ is unobstructed in $\mathbb{P}^M$. If $a = 1$, then we have that $H^1(N_{Y/\mathbb{P}^M} \otimes K_Y) = 0$ by \cite[Theorem $1.7$]{GP97} while if $a \geq 2$, the same vanishing follows from \cite[Lemma $2.7$, $(2)$]{PMR21}. and hence by Remark \ref{deformations fixing a subscheme}, the functor $F_{E, Y}$ is unobstructed. Then by Theorem \ref{main}, the result follows. \QEDA

%Consider a simple normal crossing variety $V = Y_{1t} \bigcup Y_{2t}$ as in Theorem \ref{snc to carpet} $(2)$. The sheaf $T_V^1$ is defined as follows
%$$0 \to T_V \to T_{\mathbb{P}^N}|_V \to N_{V/\mathbb{P}^N} \to T_V^1 \to 0$$
%The equisingular normal sheaf $N_{V/\mathbb{P}^N}'$ is the kernel of the map $N_{V/\mathbb{P}^N} \to T_V^1 \to 0$. Further the sheaf $T_V^1$ is line bundle on the double locus $E$, given by $T_V^1 = N_{E/Y_1} \otimes N_{E/Y_2} = \omega_Y^{-2}|_E$. We have the following lemma from \cite{CLM93}.

\begin{remark}\label{vanishings}
We record some important vanishings that we showed or used in the previous theorem. Let $Y \hookrightarrow \mathbb{P}^M$ be as in Theorem \ref{snc to carpet}. Then
\begin{enumerate}
    \item $H^1(T_{\mathbb{P}^M}|_Y) = 0$
    \item $H^1(N_{Y/\mathbb{P}^M}) = 0$
    \item $H^1(N_{Y/\mathbb{P}^M} \otimes K_Y) = 0$
\end{enumerate}
    
\end{remark}



\begin{lemma}\label{general element of the fiber of the flag Hilbert scheme}
Let $Y = \mathbb{F}_e \hookrightarrow \mathbb{P}^M$ be embedded by the complete linear series of the very ample line bundle $\mathcal{O}_Y(aC_0+bf)$ and $E \in |\omega_Y^{-1}|$ be a smooth elliptic curve. 
\begin{enumerate}
    \item[(1)] If $e \leq 1$, then a general element of $\mathcal{F}_{E,Y}$ is once again $Y_2 = \mathbb{F}_e \hookrightarrow \mathbb{P}^M$ embedded by $\mathcal{O}_{\mathbb{F}_e}(aC_0+bf)$.
    \item[(2)] If $e = 2$, then a general element of $\mathcal{F}_{E,Y}$ is $Y_2 = \mathbb{F}_0 \hookrightarrow \mathbb{P}^M$ embedded by $\mathcal{O}_{\mathbb{F}_0}(aC_0+(b-a)f) = \mathcal{O}_{\mathbb{P}^1 \times \mathbb{P}^1}(a, b-a)$.
\end{enumerate}
    
\end{lemma}

\noindent\textit{Proof.} Since the general member of any one-parameter family of deformations of $\mathbb{F}_0$ or $\mathbb{F}_1$ is once again $\mathbb{F}_0$ or $\mathbb{F}_1$ respectively, part $(1)$ follows. To show part $(2)$, we need to \sout{\textcolor{red}{need to}} show that $(E \subset Y_2 = \mathbb{F}_0 \hookrightarrow \mathbb{P}^M)$ embedded by $\mathcal{O}_Y(aC_0+(b-a)f)$ belongs to $\mathcal{F}_{E,Y}$. To see this consider a one parameter family of deformations of $\mathcal{X} \to T$, such that $\mathcal{X}_0 = \mathbb{F}_2$ and general fiber $\mathcal{X}_t = \mathbb{F}_0$. Now $E \in |-K_{\mathbb{F}_2}|$ and $H^1(-K_{\mathbb{F}_2}) = 0$. Then the elliptic curve $E$ lifts to a global divisor $\mathcal{E} \to T$ after possibly shrinking $T$. Further since $H^1(T_E) = 0$, any infinitesimal deformation of $E$ is trivial (\textbf{COMMENT:} See discussion in \url{https://math.stackexchange.com/questions/590978/first-order-and-infinitesimal-deformations}) and hence by \cite[Proposition $2.6.10$]{Ser}, after an etale base change, we can assume that $\mathcal{E} = E \times T$. Now consider the exact sequence 
$$0 \to T_Y \to T_{\mathbb{P}^M}|_Y \to N_{Y/\mathbb{P}^M} \to 0$$
Taking  cohomology we have
$$ H^0(N_{Y/\mathbb{P}^M}) \to H^1(T_Y) \to H^1(T_{\mathbb{P}^M}|_Y) \to H^1(N_{Y/\mathbb{P}^M}) \to H^2(T_Y)$$
From Remark \ref{vanishings}, $(1)$, we have that $H^1(T_{\mathbb{P}^M}|_Y)$. Hence the forgetful map between the functor $\textrm{\bf{Def}}_{Y/\mathbb{P}^M}$ of embedded deformations of $Y$ in $\mathbb{P}^M$ to the functor $\textrm{\bf{Def}}_{Y}$ of abstract deformations of $Y$ is smooth by \cite[Proposition $2.3.6$]{Ser} and consequently after possibly shrinking $T$, the family $E \times T \subset \mathcal{X} \to T$ can be lifted to a family $E \times T \subset \mathcal{X} \subset \mathbb{P}_T^M \to T$. Hence $(E \subset Y_2 = \mathbb{F}_0 \hookrightarrow \mathbb{P}^M) \in \mathcal{F}_{E,Y}$. Next we want to show that the line bundle $L_0 = \mathcal{O}_{\mathcal{X}_0}(1) = \mathcal{O}_{\mathbb{F}_2}(aC_0+bf)$ deforms to the line bundle $L_t = \mathcal{O}_{\mathcal{X}_t}(1) = \mathcal{O}_{\mathbb{P}^1 \times \mathbb{P}^1}(a, b-a)$ on the general fiber. Suppose that $L_0 = \mathcal{O}_{\mathbb{F}_2}(aC_0+bf)$ deforms to $L_t = \mathcal{O}_{\mathbb{P}^1 \times \mathbb{P}^1}(a', b')$. Then, since $L_0^2 = L_t^2$, we have that $-a^2+ab = a'b'$. Also since $h^0(L_0) = h^0(L_t) = M+1$, we have that $(b+1)(a+1)-a(a+1) = (a'+1)(b'+1) \implies a'+b' = b$. Combining the two equations, we have that $b(a'-a) = (a'+a)(a'-a)$. Suppose that $a' = a$, then we have that $b' = b-a$ and hence $\mathcal{O}_{\mathbb{P}^1 \times \mathbb{P}^1}(a', b') = \mathcal{O}_{\mathbb{P}^1 \times \mathbb{P}^1}(a, b-a)$.  On the other hand, if $b' = a'+a$, then $b' = a$ and hence $a' = b-a$. Then $\mathcal{O}_{\mathbb{P}^1 \times \mathbb{P}^1}(a', b') = \mathcal{O}_{\mathbb{P}^1 \times \mathbb{P}^1}(b-a, a) \cong \mathcal{O}_{\mathbb{P}^1 \times \mathbb{P}^1}(a, b-a)$. Hence our claim is proved. \QEDA 

\end{comment}


\begin{theorem}\label{Hirzebruch main}
Let $Y = \mathbb{F}_e \hookrightarrow \mathbb{P}^N$ be an embedding of the Hirzebruch surface $\mathbb{F}_e$, with $e \leq 2$, induced by the complete linear series of the very ample line bundle $|aC_0+bf|$ (hence $b \geq ae+1$). Let $i: \mathbb{P}^N \to \mathbb{P}^M$ be a linear embedding with $M \geq N$ and consider the composed (possibly degenerate) embedding $Y \hookrightarrow \mathbb{P}^M$. Let $E \in |\omega_Y^{-1}|$ denote a smooth elliptic curve.
\begin{itemize}
    \item[(1)] If $Y_1$ is the subvariety $Y$ of $\mathbb P^M$ and $Y_2$ is a subvariety of $\mathbb P^M$ that corresponds to a general element of $\mathcal{F}_{E,Y}$, then {the scheme theoretic intersection of $Y_1 and Y_2$ is $E$ and} {$V = Y_1 \bigcup_E Y_2$} is  smoothable inside $\mathbb{P}^M$ \sout{into}  \textcolor{red}{to a} \textbf{smooth $K3$ surface}. The Hilbert point of $V$ inside $\mathbb{P}^M$ is smooth. 
    \begin{enumerate}
        \item[(a)] If $M = N$, the subvarieties $V$ and the $K3$ surfaces are embedded by a sublinear series of codimension $$h^0((a-2)C_0+(b-e-2)f) = \displaystyle\frac{(a-1)}{2}(2b-2-ae)$$ of a very ample line bundle.
        \item[(b)] If $M = N + h^0((a-2)C_0+(b-e-2)f)$, the subvarieties $V$ and the $K3$ surfaces are embedded by a complete linear series  of a very ample line bundle.
    \end{enumerate}
     Further the closure inside the Hilbert scheme of the locus of varieties of the form $V = Y_1 \bigcup_E Y_2$ contains embedded $K3$ carpets $\widetilde{Y} \hookrightarrow \mathbb{P}^M$ which extends the embedding $Y \hookrightarrow \mathbb{P}^M$.  
    
    \smallskip
    
    \item[(2)] The locally trivial deformations of $V$ inside $\mathbb{P}^M$ are unobstructed and form a subspace in the tangent space of the Hilbert scheme at $[V]$ of codimension $h^0(\omega_Y^{-2}|_E) = 6$. Furthermore, this subspace is the tangent space at $[V]$ of an irreducible locus of the Hilbert scheme, which has codimension $h^0(\omega_Y^{-2}|_E) = 6$ and is smooth at $[V]$. The singularities of the subschemes parameterized by this locus are normal crossing singularities;  in particular, they are non-normal and semi-log-canonical. They are analytically isomorphic to $(x_1^2+x_2^2 = 0 \subset \mathbb{C}^M)$. 
   % \item[(3)] The general fiber of any one-parameter family of embedded deformations of $V$ whose image under the Kodaira-Spencer map is a  locally trivial deformation of $V$, is once again a simple normal crossing variety $V' = Y_1' \bigcup_D' Y_2'$, where both $Y_1'$ and $Y_2'$ are deformations of the subvariety $Y$ and $D'$ is a deformation of the subscheme $D$. 
    \smallskip
    
    \item[(3)] The general smooth surface of the Hilbert component containing $V$ is a 
    \begin{itemize}
        \item[(a)] prime $K3$ surface if $\textrm{gcd}(a,b) = 1$, i.e, the Picard group is generated by the hyperplane class.
        \item[(b)] a non-prime $K3$ surface where the hyperplane class is a multiple of the generator of the Picard group, the multiplicity being equal to $\textrm{gcd}(a,b)$
    \end{itemize}
\end{itemize}

\end{theorem}

\noindent\textit{Proof.} Let us prove part $(1)$. We apply Theorem \ref{intrinsic}. First of all, there exists nowhere vanishing sections inside $H^0(N_{Y/\mathbb{P}^M} \otimes K_Y)$ corresponding to embedded ribbons $\widetilde{Y}$ by \cite[Theorem $1.7$]{GP97} for $a = 1$ and  by \cite[Theorem $3.1$]{PMR21} for $a \geq 2$. By Remark \ref{vanishings}, $(2)$, $(3)$, $Y$ is unobstructed in $\mathbb{P}^M$ and the functor $F_{E, Y}$ is unobstructed. If $a = 1$, the smoothability of $\widetilde{Y}$ inside $\mathbb{P}^M$ and the smoothness of its Hilbert point follows from \cite[Corollary $2.9$]{GGP13} and \cite[Theorem $4.1$]{GP97} respectively while if $a \geq 2$, the same follows from \cite[Theorem $4.1$]{PMR21} and \cite[Theorem $5.1$]{PMR21} respectively. We remark that although the proofs of \cite[Theorem $4.1$]{PMR21} and \cite[Theorem $5.1$]{PMR21} assume that $\widetilde{Y}$ is embedded by a complete linear series, the proofs go through verbatim by embeddings of $\widetilde{Y}$ induced by an incomplete linear series. Now we prove the claim concerning the codimension of the sublinear series of the very ample line bundle that induce the embedding of the smoothing fiber. 
Consider a $K3$ double cover of $\pi:X \to Y$ branched along $B \in |\omega_Y^{-2}|$. Let $A = \pi^*(\mathcal{O}_{Y}(1))$ and consider the composition  
\[
\begin{tikzcd}
X \arrow[dr, "\varphi"] \arrow[d, "\pi"] & \\
Y \arrow[r, hook] & \mathbb{P}^N
\end{tikzcd}
\]
The smoothing fibers $X_t \xhookrightarrow{\varphi_t} \mathbb{P}^M$ are obtained as the image of an embedding $\varphi_t$ which is a deformation of $\varphi$ along a one parameter family $T$. Now note that $\varphi^* (\mathcal{O}_{\mathbb{P}^M}(1)) = \pi^*\mathcal{O}_Y(1)$. Then $A = \pi^*(\mathcal{O}_Y(aC_0+bf))$. Note that the one parameter family of smooth fibers $X_t$ that smoothes the normal crossing $V = Y_1 \bigcup_E Y_2$ are embedded inside $\mathbb{P}^M$ by some sublinear series of a line bundle $A_t$ which is a deformation of $A$. Since $X_t$ is $K3$, we have that for $i = 1,2$, $H^i(A_t) = 0$ and hence $h^0(A_t) = h^0(A)$. Hence $h^0(A_t) = h^0(A) = h^0(\pi^*(\mathcal{O}_Y(aC_0+bf))) = h^0(\mathcal{O}_Y(aC_0+bf))+h^0(\mathcal{O}_Y(aC_0+bf) \otimes \omega_Y) = N+1+h^0(\mathcal{O}_Y(aC_0+bf) \otimes \omega_Y)$. Now part (a) and (b) of $(1)$ follows by noting that $h^0(\mathcal{O}_Y(aC_0+bf) \otimes \omega_Y) = h^0(\mathcal{O}_Y((a-2)C_0+(b-e-2)f))$ whose dimension is as stated in the Theorem.

Therefore part $(1)$ follows from Theorem \ref{intrinsic}. \\
To show part $(2)$ we need to show that $H^1(N_{V/\mathbb{P}^M}') = 0$ for one subvariety of the form $V = Y_1 \bigcup_E Y_2$ where $Y_2 \in \mathcal{F}_{E,Y}$. For this, first note that for any such $V$ we have an exact sequence 
$$0 \to N_{V/\mathbb{P}^M}' \to N_{V/\mathbb{P}^M} \to T_V^1 \to 0$$
Note that by Theorem \ref{main}, we have flat family $\mathcal{Y} \to T$ over an smooth irreducible curve $T$, such that $\mathcal{Y}_0$ is an embedded $K3$ carpet $\widetilde{Y}$ on $Y$ and $\mathcal{Y}_t$ is a subvariety of the form $V = Y_1 \bigcup_E Y_2$. When $a \geq 2$, by \cite[Theorem $5.1$, equation $(5.9)$]{PMR21}, we have that $H^1(N_{\widetilde{Y}/\mathbb{P}^M}) = 0$, while for $a = 1$, the same holds for by \cite[Theorem $4.1$]{GP97}. Now since the family $\mathcal{Y} \to T$ is local complete intersection, we have that $N_{\mathcal{Y}/\mathbb{P}_T^M}$ is locally free and hence flat over $T$. Therefore by semi-continuity, $H^1(N_{\mathcal{Y}_t/\mathbb{P}^M}) = 0$, i.e, $H^1(N_{V/\mathbb{P}^M}) = 0$. The vanishing of $H^1(N_{V/\mathbb{P}^M}')$ now follows from the following two lemmas which are generalizations of \cite[Lemma $(2)$]{CLM93} and \cite[Corollary $(1)$]{CLM93} for any $a \geq 1$, whose proof works exactly as in \cite{CLM93} (note that we have already showed \textcolor{red}{in} Theorem \ref{snc to carpet} that $H^1(N_{Y/\mathbb{P}^M}) = 0$)  
\begin{lemma}\label{alternate T^1}
The sheaf $T_V^1$ is the cokernel of the inclusion $N_{Y/\mathbb{P}^M} \to N_{V/\mathbb{P}^M}|_Y$
\end{lemma}

\begin{lemma}\label{alternate smoothing}
The map of global sections from $H^0(N_{V/\mathbb{P}^M}) \to H^0(T_V^1)$ is a surjection. 
\end{lemma}

We prove part $(3)$. Let $Y = \mathbb{F}_0$ or $Y = \mathbb{F}_1$. Since the general $K3$ surface in any irreducible component of the Hilbert scheme of $\mathbb{P}^M$, parameterizing smooth $K3$ surfaces has Picard rank one, a general smoothing fiber of $V$ has Picard rank one. Once again consider the $K3$ double cover of $\pi:X \to Y$ branched along $B \in |\omega_Y^{-2}|$. Let $A = \pi^*(\mathcal{O}_{Y}(1))$ and consider the composition  
\[
\begin{tikzcd}
X \arrow[dr, "\varphi"] \arrow[d, "\pi"] & \\
Y \arrow[r, hook] & \mathbb{P}^M
\end{tikzcd}
\]
Recall the smoothing fibers $X_t \xhookrightarrow{\varphi_t} \mathbb{P}^M$ are obtained as the image of an embedding $\varphi_t$ which is a deformation of $\varphi$ along a one parameter family $T$ and are hence embedded inside $\mathbb{P}^M$ by some sublinear series of a line bundle $A_t$ which is a deformation of $A$. Since the multiplicity of a polarization remains constant on a smooth family, we have to find the multiplicity of $A$ on the double cover $X$. Now since $B$ is smooth and irreducible part $(3)$ for $Y = \mathbb{F}_0$ or $Y = \mathbb{F}_1$ follows from the fact that $\textrm{Pic}(X)$ is generated by $\pi^*\mathcal{O}_Y(C_0)$ and $\pi^*\mathcal{O}_Y(f)$ (see \cite[Proposition $6.3$]{AHL10}). Now suppose $V = Y_1 \bigcup_E Y_2$, with $Y = \mathbb{F}_2 \hookrightarrow \mathbb{P}^N$ and $\mathcal{O}_Y(1) = \mathcal{O}_{\mathbb{F}_2}(aC_0+bf)$. Then by Theorem \ref{theorem.irreducible} and Lemma \ref{general element of the fiber of the flag Hilbert scheme} part $(2)$, there is a reduced irreducible variety $\mathcal{F}_{E,Y} \times \mathcal{F}_{E,Y}$ which contains the normal crossing variety $V$ and whose general point $V' = Y_1' \bigcup_E Y_2'$ where each $Y_i' = \mathbb{F}_0 \hookrightarrow \mathbb{P}^M$ with $\mathcal{O}_{Y_i}(1) = \mathcal{O}_{\mathbb{F}_0}(aC_0+(b-a)f)$. Since the Hilbert point of $V$ and $V'$ are both smooth, they both are contained in a unique irreducible component of the Hilbert scheme. Hence by part $(3)$ applied to $V'$, the general $K3$ surface in this component of the Hilbert scheme has Picard group $\mathbb{Z}$ with the hyperplane class a multiple of the generator of the Picard group, the multiplicity being equal to $\textrm{gcd}(a,b-a) = \textrm{gcd}(a,b)$. 
%We first show that the map $\pi^*: \textrm{Pic}(Y) \to \textrm{Pic}(X)$ is an isomorphism. By \cite{Par91}, Lemma $4.2$, we have $H^1(\Omega_X^1) = H^1(\pi_*\Omega_X^1) = H^1(\Omega_Y) \bigoplus H^1(\Omega_Y(\log (B)) \otimes \omega_Y)$. Now using the exact sequence 
%$$ 0 \to \Omega_Y \otimes \omega_Y \to \Omega_Y(\log (B)) \otimes \omega_Y \to \omega_Y|_B \to 0 $$ we have 
%$$H^1(\Omega_Y \otimes K_Y) \to H^1(\Omega_Y(\log (B)) \otimes K_Y) \to H^1(K_Y|_B)$$
%Suppose that $Y = \mathbb{F}_0$ or $Y = \mathbb{F}_1$. Then the leftmost term is zero due to Nakano vanishing theorem. We have 
%$$0 \to \omega_Y^{\otimes 2} \to \omega_Y \to \omega_Y|_B \to 0$$ 
%We have $H^1(\omega_Y) = H^1(\mathcal{O}_Y) = 0$. Now applying Serre-duality, the map $H^2(\omega_Y^{\otimes 2}) \to H^2(\omega_Y)$ is the map between  and hence we see that $H^1(\omega_Y|_B) = 0$. This implies $H^1(\Omega_Y(\log (B)) \otimes K_Y) = 0$. Now since $H^1(\mathcal{O}_Y) = H^1(\mathcal{O}_X) = 0$, the pullback map between $\pi^*: \textrm{Pic}(Y)  \to \textrm{Pic}(X) $ is induced by the map 
%$$H^1(\Omega_Y) \to H^1(\Omega_Y) \bigoplus H^1(\Omega_Y(\log (B)) \otimes K_Y)$$
%Therefore $H^1(\Omega_Y(\log (B)) \otimes K_Y) = 0$ implies $\pi^*$ is an isomorphism


\QEDA
\color{black}


\begin{comment}
\subsection{Syzygy points of union of scrolls and their semistability}

We refer to \cite{Fed18} for the relevant defintions of syzygy points and their semistabilty.  

\begin{corollary}
    \begin{itemize}
        
        \item[(1)] For each $g \geq 3$ and $g = 2m+1$, a general smoothable simple normal crossing union of balanced scrolls $V = Y_1 \bigcup_E Y_2 \subset \mathbb{P}^g$ meeting along an anticanonical elliptic curve $E$ as constructed in Theorem \ref{Hirzebruch main}, where $Y_i = \mathbb{F}_0 \hookrightarrow \mathbb{P}^g$ is embedded by the complete linear series of $|C_0+mf|$, has well-defined syzygy points $\textrm{Syz}_{0,q}(V)$ for $q \geq 2$ and $\textrm{Syz}_{p,2}(V)$ for $1 \leq p \leq m-1$, which are semistable with respect to the $SL(g+1)$ action.

        \item[(2)] For each $g \geq 3$ and $g = 2m+1$, a general canonical binary curve $D = C_1 \bigcup C_2$ of genus $g$, where $C_i$'s are rational normal curves of degree $g-1$ intersecting at $g+1$ points has well-defined syzygy points $\textrm{Syz}_{0,q}(D)$ and $\textrm{Syz}_{1,2}(D)$ which are semistable with respect to the $SL(g)$ action.
        
        
    \end{itemize}
    
\end{corollary}

\noindent\textit{Proof $(1)$.} By Theorem \ref{Hirzebruch main} (see also Theorem \ref{snc to carpet}) there exists a flat family of subschemes $\sigma: \mathcal{Y} \to T$ of $\mathbb{P}^g_T$ over a smooth irreducible algebraic curve $T$ such that 
\begin{itemize}
    \item[(1)]  $\mathcal{Y}_0$ is an embedded $K3$ carpet $\widetilde{Y} \hookrightarrow \mathbb{P}^g$ extending an embedding $Y = \mathbb{F}_0 \hookrightarrow \mathbb{P}^g$.  
    \item[(2)] $\mathcal{Y}_t = Y_{1t} \bigcup_E Y_{2t}$, where for each $t$ and $i = 1,2$, $Y_{it} = \mathbb{F}_0 \hookrightarrow \mathbb{P}^g$ is embedded by the complete linear series of $|C_0+mf|$ and $E$ is a smooth elliptic curve, anticanonical in each $Y_{it}$.
    \item[(3)] The varieties $\mathcal{Y}_t$ are smoothable in $\mathbb{P}^{g}$ into smooth $K3$ surfaces.
\end{itemize}

By Theorem \ref{Hirzebruch main}, part $(1)$, (b), $\mathcal{Y}_0 \hookrightarrow \mathbb{P}^g$ is embedded by a complete linear series. It is also known that a $K3$ carpet on $\mathbb{F}_0$ embedded as a scroll is projectively normal (see \cite{BE95},  \cite{Deo18}). (\textcolor{red}{I do not understand meaning of "projectively normal" when let alone the cone, the scheme itself is not normal". This is a grotesque abuse of term, whoever may have used it. One can of course talk of the multiplication map.}) Consider the line bundle $\mathcal{L} = \mathcal{O}_{\mathcal{Y}}(1)$ and the cokernel $\mathcal{G}$ of the map 
$$ \sigma_*(\mathcal{O}_{\mathbb{P}^g_T}(n)) \to  \sigma_*(\mathcal{O}_{\mathcal{Y}}(n)) \to \mathcal{G} \to 0$$
$\mathcal{G}$ is a coherent sheaf on $T$, 
\textcolor{red}{Is it better if we write $\mathcal{G}$ explicitly as $R^1\sigma_*(\mathcal{I}(n))$}?  
and by cohomology and base change, $\mathcal{G}_t = \textrm{coker}(\mathcal{O}_{\mathbb{P}^g}(n) \to  \mathcal{O}_{\mathcal{Y}_t}(n))$. Now since $\mathcal{G}_0 = (0)$ we have that $\mathcal{G}_t = (0)$ for a general $t$ and hence $\widetilde{Y}_t \hookrightarrow \mathbb{P}^g$ is embedded by a complete linear series and is projectively normal. Therefore the syzygy points $\textrm{Syz}_{p,q}(\mathcal{Y}_t)$ are well-defined as soon as $K_{p,q}(\mathcal{Y}_t) = 0$.  Now it has been shown in \cite{Deo18} that the syzygy points $K_{p,2}(\mathcal{Y}_0) = 0$ for $1 \leq p \leq m-1$ are well-defined. Let $\mathcal{K}_{p,2}$ denote the cohomology in the middle of the sequence 
$$\bigwedge^{p+1} \sigma_* \mathcal{L} \otimes \sigma_* \mathcal{L} \to \bigwedge^p \sigma_* \mathcal{L} \otimes \sigma_* \mathcal{L}^{\otimes 2} \to \bigwedge^{p-1} \sigma_* \mathcal{L} \otimes \sigma_* \mathcal{L}^{\otimes 3}$$
As before $\mathcal{K}_{p,2}$ is a coherent sheaf (since kernel and cokernel of coherent sheaves are coherent) and by cohomology and base change we have that the fiber of $\mathcal{K}_{p,2}$ at $t \in T$ is $K_{p,2}(\mathcal{Y}_t)$, we conclude that $K_{p,2}(\mathcal{Y}_t) = 0$ for for $1 \leq p \leq m-1$. Now by openness of semistability, the result follows from the proof semistability of syzygy points $\textrm{Syz}_{0,q}(\mathcal{Y}_0)$ for $q \geq 2$ and $\textrm{Syz}_{p,2}(\mathcal{Y}_0)$ for $1 \leq p \leq m-1$ in \cite{Fed18}, proof of Theorem $4.1$. 

\noindent\textit{Proof.} Taking a general hyperplane section of the family obtained in part $(1)$, we get a flat family of subschemes $\sigma: \mathcal{D} \to T$ of $\mathbb{P}^{g-1}_T$ over a smooth irreducible algebraic curve $T$ such that 
\begin{itemize}
    \item[(1)]  $\mathcal{D}_0$ is a canonical ribbon $\widetilde{D} \hookrightarrow \mathbb{P}^{g-1}$ extending an embedding $C = \mathbb{P}^1 \hookrightarrow \mathbb{P}^{g-1}$.  
    \item[(2)] $\mathcal{D}_t = C_{1t} \bigcup C_{2t}$, where for each $t$ and $i = 1,2$, $C_{it}$ is a rational normal curve of degree $g-1$, $C_{1t} \bigcap C_{2t}$ is a reduced scheme of length $g+1$ and $\mathcal{D}_t$ is canonically embedded in $\mathbb{P}^{g-1}$.
    \item[(3)] The varieties $\mathcal{D}_t$ are smoothable in $\mathbb{P}^{g-1}$ into smooth canonical curves.
\end{itemize}

Now once again the well-definedness of the syzygy points follow by \cite{Deo18} while the semistability 
$\textrm{Syz}_{0,q}(\mathcal{D}_t)$ and $\textrm{Syz}_{1,2}(\mathcal{D}_t)$ follows from the corresponding semistability of $\textrm{Syz}_{0,q}(\mathcal{D}_0)$ and $\textrm{Syz}_{1,2}(\mathcal{D}_0)$ (see \cite{AFS13}, \cite{FJ13} \cite{DFS16}).
\QEDA



\subsection{Explicit examples of degenerations of $K3$ surfaces of genus $g$ for some low values of $g$}

We fix a $g \geq 3$ and consider the Hilbert scheme of $K3$ surfaces $X \hookrightarrow \mathbb{P}^g$ embedded by the complete linear series of a line bundle $L$, with $L^2 = 2g-2$. We list below all the distinct codimension six locus consisting of snc degenerations into a union of two embedded Hirzebruch surfaces, as obtained from Theorem \ref{Hirzebruch main} for some low values of $g$. We compare our results with \cite{CLM93} and point out the new aspects of our results. 

\begin{example}
    \begin{enumerate}
        \item For $g = 3$, we get a locus of snc varieties of the form $V = Y_1 \bigcup_E Y_2 \hookrightarrow \mathbb{P}^3$, where each $Y_i = \mathbb{F}_0 \hookrightarrow \mathbb{P}^3$ is embedded by $\mathcal{O}_{\mathbb{F}_0}(C_0+f)$. The locus is contained in the unique component of the Hilbert scheme parameterizing prime $K3$ surfaces of genus $3$. \textbf{The new result is that the Hilbert points of general such $V$ are semistable.}
        \smallskip
        \item For $g = 4$, we get a locus of snc varieties of the form $V = Y_1 \bigcup_E Y_2 \hookrightarrow \mathbb{P}^4$, where each $Y_i = \mathbb{F}_1 \hookrightarrow \mathbb{P}^4$ is embedded by $\mathcal{O}_{\mathbb{F}_1}(C_0+2f)$. The locus is contained in the unique component of the Hilbert scheme parameterizing prime $K3$ surfaces of genus $4$. 
        \smallskip
        \item For $g = 5$, we get a locus of snc varieties of the form $V = Y_1 \bigcup_E Y_2 \hookrightarrow \mathbb{P}^5$, where each $Y_i = \mathbb{F}_0 \hookrightarrow \mathbb{P}^5$ is embedded by $\mathcal{O}_{\mathbb{F}_0}(C_0+2f)$. The locus is contained in the unique component of the Hilbert scheme parameterizing prime $K3$ surfaces of genus $5$. \textbf{The new result is that the Hilbert points of general such $V$ are semistable.} There is a sublocus consisting of varieties of the form $V = Y_1 \bigcup_E Y_2 \hookrightarrow \mathbb{P}^5$, where $Y_1 = \mathbb{F}_0 \hookrightarrow \mathbb{P}^5$ is embedded by $\mathcal{O}_{\mathbb{F}_0}(C_0+2f)$ and $Y_2 = \mathbb{F}_2 \hookrightarrow \mathbb{P}^5$ is embedded by $\mathcal{O}_{\mathbb{F}_2}(C_0+3f)$. 
        \smallskip
        \item For $g = 6$, we get a locus of snc varieties of the form $V = Y_1 \bigcup_E Y_2 \hookrightarrow \mathbb{P}^6$, where each $Y_i = \mathbb{F}_1 \hookrightarrow \mathbb{P}^6$ is embedded by $\mathcal{O}_{\mathbb{F}_1}(C_0+3f)$. The locus is contained in the unique component of the Hilbert scheme parameterizing prime $K3$ surfaces of genus $6$
        \smallskip
        \item For $g = 7$, we get a locus of snc varieties of the form $V = Y_1 \bigcup_E Y_2 \hookrightarrow \mathbb{P}^7$, where each $Y_i = \mathbb{F}_1 \hookrightarrow \mathbb{P}^7$ is embedded by $\mathcal{O}_{\mathbb{F}_0}(C_0+3f)$. The locus is contained in the unique component of the Hilbert scheme parameterizing prime $K3$ surfaces of genus $7$. \textbf{The new result is that the Hilbert points of general such $V$ are semistable.}
        \smallskip
        \item For $g = 8$, we get a locus of snc varieties of the form $V = Y_1 \bigcup_E Y_2 \hookrightarrow \mathbb{P}^7$, where each $Y_i = \mathbb{F}_1 \hookrightarrow \mathbb{P}^8$ is embedded by $\mathcal{O}_{\mathbb{F}_1}(C_0+4f)$. The locus is contained in the unique component of the Hilbert scheme parameterizing prime $K3$ surfaces of genus $8$. 
        \smallskip
        \item For $g = 9$, we get three different loci 
        \begin{enumerate}
            \item a locus of snc varieties of the form $V = Y_1 \bigcup_E Y_2 \hookrightarrow \mathbb{P}^9$, where each $Y_i = \mathbb{F}_0 \hookrightarrow \mathbb{P}^9$ is embedded by $\mathcal{O}_{\mathbb{F}_1}(C_0+4f)$. The locus is contained in the unique component of the Hilbert scheme parameterizing prime $K3$ surfaces of genus $9$. \textbf{The new result is that the Hilbert points of general such $V$ are semistable.} There is a sublocus consisting of varieties of the form $V = Y_1 \bigcup_E Y_2 \hookrightarrow \mathbb{P}^9$, where $Y_1 = \mathbb{F}_0 \hookrightarrow \mathbb{P}^8$ is embedded by $\mathcal{O}_{\mathbb{F}_0}(C_0+4f)$ and $Y_2 = \mathbb{F}_2 \hookrightarrow \mathbb{P}^8$ is embedded by $\mathcal{O}_{\mathbb{F}_2}(C_0+5f)$.  
            \item a locus of snc varieties of the form $V = Y_1 \bigcup_E Y_2 \hookrightarrow \mathbb{P}^9$, where each $Y_i = \mathbb{F}_1 \hookrightarrow \mathbb{P}^8$ is embedded by $\mathcal{O}_{\mathbb{F}_1}(2C_0+3f)$. The locus is contained in the unique component of the Hilbert scheme parameterizing prime $K3$ surfaces of genus $9$. \textbf{The existence of this locus is new and in fact by Theorem \ref{Hirzebruch main}, we show that the degeneration actually happens in $\mathbb{P}^{8}$.}
            \item a locus of snc varieties of the form $V = Y_1 \bigcup_E Y_2 \hookrightarrow \mathbb{P}^9$, where each $Y_i = \mathbb{F}_0 \hookrightarrow \mathbb{P}^9$ is embedded by $\mathcal{O}_{\mathbb{F}_0}(2C_0+2f)$. The locus is contained in the component of the Hilbert scheme parameterizing non-prime $K3$ surfaces of genus $9$ where the hyperplane bundle is twice the generator of the Picard group. \textbf{This degeneration can be obtained by taking the second veronese embedding of the degeneration in item $(1)$.} \textbf{But the new result is that by Theorem \ref{Hirzebruch main}, we show that this degeneration actually happens in $\mathbb{P}^8$.}
        \end{enumerate} 
        \smallskip
        \item For $g = 13$, we get three different loci 
        \begin{enumerate}
            \item a locus of snc varieties of the form $V = Y_1 \bigcup_E Y_2 \hookrightarrow \mathbb{P}^{13}$, where each $Y_i = \mathbb{F}_0 \hookrightarrow \mathbb{P}^{13}$ is embedded by $\mathcal{O}_{\mathbb{F}_0}(C_0+6f)$. The locus is contained in the unique component of the Hilbert scheme parameterizing prime $K3$ surfaces of genus $13$. \textbf{The new result is that the Hilbert points of general such $V$ are semistable.}
            \smallskip
            \item a locus of snc varieties of the form $V = Y_1 \bigcup_E Y_2 \hookrightarrow \mathbb{P}^{13}$, where each $Y_i = \mathbb{F}_0 \hookrightarrow \mathbb{P}^{11}$ is embedded by $\mathcal{O}_{\mathbb{F}_0}(2C_0+3f)$. The locus is contained in the unique component of the Hilbert scheme parameterizing prime $K3$ surfaces of genus $13$. \textbf{The existence of this locus is new and in fact by Theorem \ref{Hirzebruch main}, we show that this degeneration actually happens in $\mathbb{P}^{11}$.} There is a sublocus consisting of varieties of the form $V = Y_1 \bigcup_E Y_2 \hookrightarrow \mathbb{P}^{13}$, where $Y_1 = \mathbb{F}_0 \hookrightarrow \mathbb{P}^{11}$ is embedded by $\mathcal{O}_{\mathbb{F}_0}(2C_0+3f)$ and $Y_2 = \mathbb{F}_2 \hookrightarrow \mathbb{P}^{11}$ is embedded by $\mathcal{O}_{\mathbb{F}_2}(2C_0+5f)$.
            \smallskip
            \item a locus of snc varieties of the form $V = Y_1 \bigcup_E Y_2 \hookrightarrow \mathbb{P}^{13}$, where each $Y_i = \mathbb{F}_1 \hookrightarrow \mathbb{P}^{11}$ is embedded by $\mathcal{O}_{\mathbb{F}_1}(2C_0+4f)$. The locus is contained in the component of the Hilbert scheme parameterizing non-prime $K3$ surfaces of genus $9$ where the hyperplane bundle is twice the generator of the Picard group. \textbf{This degeneration can be obtained by taking the second veronese embedding of the degeneration in item $(2)$.} \textbf{But the new result is that by Theorem \ref{Hirzebruch main}, we show that this degeneration actually happens in $\mathbb{P}^{11}$.}
        \end{enumerate} 
        
        
    \end{enumerate}
\end{example}



\smallskip
\end{comment}

\section{Extendability of $K3$ surfaces}

%\begin{theorem}\label{ZL for snc}
% Let $V = Y_1 \bigcup_E Y_2 \hookrightarrow \mathbb{P}^M$ be a simple normal crossing variety of dimension $d > 1$, embedded by the complete linear series of a very ample line bundle $\widetilde{H}$, where $Y_1 = \mathbb{F}_e \hookrightarrow \mathbb{P}^N$, $e \leq 2$ is embedded by the complete linear series of a very ample line bundle $|aC_0+bf|$, $E \in |-K_{Y_1}|$, $Y_2$ is a general element of $\mathcal{F}_{E,Y}$ and $M = N + h^0((a-2)C_0+(b-e-2)f)$ (i.e, we are in the situation of Theorem \ref{Hirzebruch main}, $(1)$ (b). Let $H = \widetilde{H}|_{Y_1}$ and assume $H^1(-H+K_{Y_1}) = 0$. Let



% $$\beta = h^1(T_{Y_1}(-H+K_{Y_1})) + h^1(T_{Y_1}(-H)) - h^0(T_{Y_1}(-H)) +  h^1(-H-K_{Y_1}) + h^0(-H-2K_{Y_1})$$ 

\begin{theorem}\label{ZL for K3 carpets}
 Let $\widetilde{Y} \hookrightarrow \mathbb{P}^M$ be a $K3$ carpet embedded by the complete linear series of a very ample line bundle $\widetilde{H}$ such that the embedding induced on the reduced part is $Y = \mathbb{F}_e \hookrightarrow \mathbb{P}^N$ embedded by the complete linear series of a very ample line bundle $H = aC_0+bf$. Assume $H^1(-H+K_Y) = 0$. 
 
 \begin{enumerate}
 \item Let $$\beta = h^1(T_Y(-H+K_Y)) + h^1(T_Y(-H)) - h^0(T_Y(-H)) +  h^1(-H-K_Y) + h^0(-H-2K_Y)$$ 
Then 
$$\alpha(\widetilde{Y}) = h^0(N_{\widetilde{Y}/\mathbb{P}^M}(-\widetilde{H}))-M-1 \leq  \beta$$ 
In particular if $\beta = 0$, then $\alpha(\widetilde{Y}) \leq 0$. Consequently, the general smooth $K3$ surfaces in the Hilbert component containing $\widetilde{Y}$ is not extendable .

\item Let $a = 2$, $\widetilde{Y}$ a general embedded $K3$ carpet represented by a general no-where vanishing section in $H^0(N_{Y/\mathbb{P}^M} \otimes K_Y)$ and set
%and $$\gamma = h^0(N_{Y/\mathbb{P}^M}(-H+K_Y)+h^0(N_{Y/\mathbb{P}^M}(-H)+h^0(-H-2K_Y)$$
\begin{align*}
    \gamma & =  h^0(N_{Y/\mathbb{P}^M}(-H+K_Y)+h^0(N_{Y/\mathbb{P}^M}(-H)+h^0(-H-2K_Y)-M-1 \\
    & \leq h^1(T_Y(-H+K_Y))+h^0(N_{Y/\mathbb{P}^M}(-H)+h^0(-H-2K_Y)-M-1
\end{align*}

Then 
$$\alpha(\widetilde{Y}) = h^0(N_{\widetilde{Y}/\mathbb{P}^M}(-\widetilde{H}))-M-1 \leq  \gamma$$ 
In particular if $\gamma = 0$, then $\alpha(\widetilde{Y}) \leq 0$. Consequently, the general smooth $K3$ surfaces in the Hilbert component containing $\widetilde{Y}$ is not extendable.


\end{enumerate}
\end{theorem}

\begin{proof}
    We prove part $(2)$. We have the following exact sequences (see \cite[Lemma $4.1$]{GGP08})
\begin{equation}\label{1}
    0 \to N_{\widetilde{Y}/\mathbb{P}^M} (K_Y) \to N_{\widetilde{Y}/\mathbb{P}^M} \to N_{\widetilde{Y}/\mathbb{P}^M} \otimes \mathcal{O}_Y \to 0
\end{equation}
\begin{equation}\label{2}
    0 \to \mathcal{H}om(I_{\widetilde{Y}}/I_{Y}^2, \mathcal{O}_Y) \to N_{\widetilde{Y}/\mathbb{P}^M} \otimes \mathcal{O}_Y \to \mathcal{O}_Y(-2K_Y) \to 0 
\end{equation}
\begin{equation}\label{3}
    0 \to \mathcal{O}_Y(-K_Y) \to N_{Y/\mathbb{P}^M} \to \mathcal{H}om(I_{\widetilde{Y}}/I_{Y}^2, \mathcal{O}_Y) \to 0
\end{equation}


Tensoring equation \ref{1} by $\mathcal{O}_{\widetilde{Y}}(-\widetilde{H})$, we have that 

\begin{equation}\label{4}
 h^0(N_{\widetilde{Y}/\mathbb{P}^M}(-\widetilde{H})) \leq h^0(N_{\widetilde{Y}/\mathbb{P}^M}(-H+K_Y)) + h^0(N_{\widetilde{Y}/\mathbb{P}^M}(-H))
\end{equation}

We compute $h^0(N_{\widetilde{Y}/\mathbb{P}^M}(-H+K_Y))$. Tensoring equation \ref{2} by $\mathcal{O}_Y(-H+K_Y)$ we have 

\begin{equation*}
    0 \to \mathcal{H}om(I_{\widetilde{Y}}/I_{Y}^2, \mathcal{O}_Y)(-H+K_Y) \to N_{\widetilde{Y}/\mathbb{P}^M}(-H+K_Y) \to \mathcal{O}_Y(-H-K_Y) \to 0 
\end{equation*}

Hence 

\begin{equation}\label{6}
    h^0(N_{\widetilde{Y}/\mathbb{P}^M}(-H+K_Y)) \leq h^0(\mathcal{H}om(I_{\widetilde{Y}}/I_{Y}^2, \mathcal{O}_Y)(-H+K_Y)) + h^0(\mathcal{O}_Y(-H-K_Y))
\end{equation}

Tensoring equation \ref{3}, by $\mathcal{O}_Y(-H+K_Y)$ we have 

\begin{equation*}
    0 \to \mathcal{O}_Y(-H) \to N_{Y/\mathbb{P}^M}(-H+K_Y) \to \mathcal{H}om(I_{\widetilde{Y}}/I_{Y}^2, \mathcal{O}_Y)(-H+K_Y) \to 0
\end{equation*}

Since $H^1(-H) = 0$, by Kodaira vanishing theorem, and $h^0(-H) = 0$, we have that 

\begin{equation}\label{7}
   h^0(\mathcal{H}om(I_{\widetilde{Y}}/I_{Y}^2, \mathcal{O}_Y)(-H+K_Y)) = h^0(N_{Y/\mathbb{P}^M}(-H+K_Y))  
\end{equation}

We now compute $h^0(N_{\widetilde{Y}/\mathbb{P}^M}(-H))$. Tensoring equation \ref{2} by $\mathcal{O}_Y(-H)$ we have 

\begin{equation*}
    0 \to \mathcal{H}om(I_{\widetilde{Y}}/I_{Y}^2, \mathcal{O}_Y)(-H) \to N_{\widetilde{Y}/\mathbb{P}^M}(-H) \to \mathcal{O}_Y(-H-2K_Y) \to 0 
\end{equation*}

Hence 

\begin{equation}\label{7}
    h^0(N_{\widetilde{Y}/\mathbb{P}^M}(-H)) \leq h^0(\mathcal{H}om(I_{\widetilde{Y}}/I_{Y}^2, \mathcal{O}_Y)(-H)) + h^0(\mathcal{O}_Y(-H-2K_Y))
\end{equation}

Tensoring equation \ref{3}, by $\mathcal{O}_Y(-H)$ we have 

\begin{equation*}
    0 \to \mathcal{O}_Y(-H-K_Y) \to N_{Y/\mathbb{P}^M}(-H) \to \mathcal{H}om(I_{\widetilde{Y}}/I_{Y}^2, \mathcal{O}_Y)(-H) \to 0
\end{equation*}
Now we claim that 

\begin{claim}\label{sp
lits}
    For a general ribbon $\widetilde{Y}$, the induced map $H^0(N_{Y/\mathbb{P}^M}(-H)) \to H^0(\mathcal{H}om(I_{\widetilde{Y}}/I_{Y}^2, \mathcal{O}_Y)(-H))$ is surjective.
\end{claim}

Granting the claim for now, we have that

\begin{equation}\label{7}
   h^0(\mathcal{H}om(I_{\widetilde{Y}}/I_{Y}^2, \mathcal{O}_Y)(-H)) \leq h^0(N_{Y/\mathbb{P}^M}(-H)) - h^0(-H-K_Y) 
\end{equation}



Hence we have that 

\begin{equation*}
 h^0(N_{\widetilde{Y}/\mathbb{P}^M}(-\widetilde{H})) \leq h^0(-H-K_Y) + h^0(N_{Y/\mathbb{P}^M}(-H+K_Y)) + h^0(N_{Y/\mathbb{P}^M}(-H)) - h^0(-H-K_Y) + h^0(-H-2K_Y)
\end{equation*}

\begin{equation}\label{8}
h^0(N_{\widetilde{Y}/\mathbb{P}^M}(-\widetilde{H})) \leq  h^0(N_{Y/\mathbb{P}^M}(-H+K_Y)) + h^0(N_{Y/\mathbb{P}^M}(-H)) + h^0(-H-2K_Y)
\end{equation}

We now compute $h^0(N_{Y/\mathbb{P}^M}(-H+K_Y))$. 
We have

\begin{equation}\label{9}
   0 \to T_Y \to T_{\mathbb{P}^M}|_Y \to N_{Y/\mathbb{P}^M} \to 0 
\end{equation}




and the Euler exact sequence restricted to $Y$

\begin{equation}\label{11}
   0 \to \mathcal{O}_Y \to \mathcal{O}_Y(H)^{M+1} \to T_{\mathbb{P}^M}|_Y \to 0 
\end{equation}





Tensoring equation \ref{11} by $\mathcal{O}_Y(-H+K_Y)$, we have 

\begin{equation*}
     H^0(\mathcal{O}_Y(K_Y))^{M+1} \to H^0(T_{\mathbb{P}^M}|_Y(-H+K_Y)) \to H^1(\mathcal{O}_Y(-H+K_Y)) 
\end{equation*}

We have that $H^0(\mathcal{O}_Y(K_Y)) = 0$ and by assumption $H^1(\mathcal{O}_Y(-H+K_Y)) = 0$ and hence $H^0(T_{\mathbb{P}^M}|_Y(-H+K_Y)) = 0$. 

%Similarly we have from the same exact sequence

%\begin{equation*}
 %   H^1(\mathcal{O}_Y(K_Y))^{N+1} \to H^1(T_{\mathbb{P}^M}|_Y(-H+K_Y)) \to H^2(\mathcal{O}_Y(-H+K_Y))
%\end{equation*}


Tensoring equation \ref{9} by $\mathcal{O}_Y(-H+K_Y)$, we have 

\begin{equation*}
    H^0(T_{\mathbb{P}^M}|_Y(-H+K_Y)) \to H^0(N_{Y/\mathbb{P}^M}(-H+K_Y)) \to H^1(T_Y(-H+K_Y))
\end{equation*}

Since $H^0(T_{\mathbb{P}^M}|_Y(-H+K_Y)) = 0$, 

\begin{equation}\label{14}
    h^0(N_{Y/\mathbb{P}^M}(-H+K_Y)) \leq h^1(T_Y(-H+K_Y))
\end{equation}

Plugging everything back into equation \ref{8}, we have

\begin{equation}\label{15}
h^0(N_{\widetilde{Y}/\mathbb{P}^M}(-\widetilde{H})) \leq  M+1 + h^1(T_Y(-H+K_Y)) + h^0(N_{Y/\mathbb{P}^M}(-H)) + h^0(-H-2K_Y)
\end{equation}



\end{proof}


We now prove the claim.

\smallskip

\textit{Proof of Claim \ref{splits}.} Consider the exact sequence given by twisting equation \ref{3} by $-H$. 
\begin{equation}\label{16}
0 \to -K_Y-H \to N_{Y/\mathbb{P}^M}(-H) \to \mathcal{H}om(I_{\widetilde{Y}}/I_{Y}^2, \mathcal{O}_Y)(-H) \to 0
\end{equation}

Let $p: Y \to \mathbb{P}^1$ denote the smooth fibration. Now note that $-K_Y-H = -(b-e-2)f$ and hence $R^1p_*(-H-K_Y) = 0$. Hence we have the following exact sequence by pushing forward the exact sequence \ref{16}

\begin{equation}\label{17}
    0 \to p_*(-H-K_Y) \to p_*(N_{Y/\mathbb{P}^M}(-H)) \to p_*(\mathcal{H}om(I_{\widetilde{Y}}/I_{Y}^2, \mathcal{O}_Y)(-H)) \to 0
\end{equation}
    
It is enough to show that the exact sequence \ref{17} is split for a general ribbon $\widetilde{Y}$. To see this we first show that there is a surjection from $p_*(N_{Y/\mathbb{P}^M}(-H))$ to $p_*(-H-K_Y)$. To see this consider the exact sequence \ref{9} twisted by $-H$ to get 

\begin{equation}\label{18}
   0 \to T_Y(-H) \to T_{\mathbb{P}^M}|_Y(-H) \to N_{Y/\mathbb{P}^M}(-H) \to 0 
\end{equation}

We show that

\begin{equation*}
   R^1p_*(T_{\mathbb{P}^M}|_Y(-H)) = 0, p_*(T_Y(-H)) = \mathcal{O}_{\mathbb{P}^1}(-(b-e)) \operatorname{and} R^1p_*(T_Y(-H)) =  \mathcal{O}_{\mathbb{P}^1}(-(b-e-2))  
\end{equation*}

To see the first vanishing once again note that from the Euler exact sequence pulled back to $Y$ and twisted by $-H$, we have 

\begin{equation}\label{24}
    0 \to \mathcal{O}_Y(-H) \to \mathcal{O}_Y^{\otimes M+1} \to T_{\mathbb{P}^M}|_Y(-H) \to 0
\end{equation}

Since $R^1p_*(\mathcal{O}_Y) = R^2(\mathcal{O}_Y(-H)) = 0$, we have that $R^1p_*(T_{\mathbb{P}^M}|_Y(-H)) = 0$. Now we turn to $R^ip_*(T_Y(-H))$. Consider the exact sequence of horizontal and vertical tangent bundles on $Y$ twisted by $-H$
\begin{equation}\label{19}
    0 \to T_{Y/\mathbb{P}^1}(-H) \to T_Y(-H) \to p^*T_{\mathbb{P}^1}(-H) \to 0
\end{equation}

Using the relative Euler exact sequence we have that 
$T_{Y/\mathbb{P}^1} = 2C_0+ef$, so that $T_{Y/\mathbb{P}^1}(-H) = -(b-e)f$. Therefore $p_*(T_{Y/\mathbb{P}^1}(-H)) = \mathcal{O}_{\mathbb{P}^1}(-(b-e))$ and $R^ip_*(T_{Y/\mathbb{P}^1}(-H)) = R^ip_*(-(b-e)f) = 0$ for $i = 1,2$. Now $p^*T_{\mathbb{P}^1}(-H) = -2C_0-(b-2)f$. Therefore, $p_*p^*T_{\mathbb{P}^1}(-H) = p_*(-2C_0-(b-2)f) = \mathcal{O}_{\mathbb{P}^1}(-(b-2)) \otimes p_*(-2C_0) = 0$ and

\begin{align*}\label{20}
  R^1p_*(p^*T_{\mathbb{P}^1}(-H)) & = R^1p_*(-2C_0-(b-2)f) \\
  & = \mathcal{O}_{\mathbb{P}^1}(-(b-2)) \otimes R^1p_*(-2C_0) \\
  & = \mathcal{O}_{\mathbb{P}^1}(-(b-2)) \otimes (p_*(\mathcal{O}_Y))^* \otimes (\bigwedge^2(\mathcal{O}_{\mathbb{P}^1} \oplus \mathcal{O}_{\mathbb{P}^1}(-e)))^* \\ 
  & = \mathcal{O}_{\mathbb{P}^1}(-(b-e-2))
\end{align*}

Pushing forward the sequence \ref{19}, we have that 

\begin{equation}\label{21}
0 \to p_*(T_{Y/\mathbb{P}^1}(-H)) = \mathcal{O}_{\mathbb{P}^1}(-(b-e)) \to p_*T_Y(-H) \to 0    
\end{equation}

and 

\begin{equation}\label{22}
   0 \to R^1p_*T_Y(-H) \to \mathcal{O}_{\mathbb{P}^1}(-(b-e-2)) \to 0 
\end{equation}

%Now pushing forward the sequence \ref{19}, we have that $R^1p_*(T_Y(-H)) = \mathcal{O}_{\mathbb{P}^1}(-(b-e-2)) $ and we have an isomorphism $p_*(T_Y(-H)) = \mathcal{O}_{\mathbb{P}^1}(-(b-e))$. 
Now pushing forward the exact sequence \ref{18}, we have that 

\begin{multline}\label{23}
0 \to \mathcal{O}_{\mathbb{P}^1}(-(b-e)) = p_*T_Y(-H) = p_*(T_{Y/\mathbb{P}^1}(-H)) \to p_*(T_{\mathbb{P}^M}|_Y(-H)) \to p_*(N_{Y/\mathbb{P}^M}(-H)) \\ \xrightarrow[]{\lambda} \mathcal{O}_{\mathbb{P}^1}(-(b-e-2)) = p_*(-H-K_Y) \to 0
\end{multline}

We now end the proof of the claim by showing that for a general ribbon $\widetilde{Y}$, which is defined by an injection $0 \to -K_Y \to N_{Y/\mathbb{P}^M}$, the map $\lambda$ is a section to the injective map induced by the pushforward after twisting by $-H$. 

$$p_*(-H-K_Y) \to p_*(N_{Y/\mathbb{P}^M}(-H))$$

So the question is can we choose a ribbon or equivalently, an injection $0 \to -K_Y \to N_{Y/\mathbb{P}^M}$, such that the induced map $\beta$,  $$0 \to p_*(-H-K_Y) = \mathcal{O}_{\mathbb{P}^1}(-(b-e-2)) \xrightarrow[]{\beta} p_*(N_{Y/\mathbb{P}^M}(-H))$$

which sits in the following diagram 


\begin{tikzcd}\label{section}
  &&& 0 \arrow[d] && \\
  &&& p_*(-H-K_Y) \arrow[d, "\beta"] && \\
  0 \arrow[r] & \mathcal{O}_{\mathbb{P}^1}(-(b-e)) \arrow[r] &  p_*(T_{\mathbb{P}^M}|_Y(-H)) \arrow[r, "\delta"] & p_*(N_{Y/\mathbb{P}^M}(-H)) \arrow[r, "\lambda"] & p_*(-H-K_Y) \arrow[r] & 0  
\end{tikzcd}

\smallskip 

satisfies $\lambda \circ \beta = \operatorname{id}$. It is enough to choose $\beta$ such that $\lambda \circ \beta$  is an isomorphism and since $p_*(-H-K_Y) = \mathcal{O}_{\mathbb{P}^1}(-(b-e-2))$ is a line bundle, it is enough to choose $\beta$ such that $\lambda \circ \beta$ is non-zero. This happens in particular if $\operatorname{Im}(\beta)$ is not contained in $\operatorname{Ker}(\lambda)$. So we need to show that $\beta$ does not factor as 

\[
\begin{tikzcd}
   & & p_*(-H-K_Y) \arrow[d, "\beta"] \arrow[dl] \\
   0 \arrow[r] & \operatorname{Ker}(\lambda) \arrow[r] & p_*(N_{Y/\mathbb{P}^M}(-H))
\end{tikzcd}
\]

It is enough to show that 

\begin{equation}\label{32}
    \operatorname{Hom}(p_*(-H-K_Y), \operatorname{Ker}(\lambda)) \to \operatorname{Hom}(p_*(-H-K_Y), p_*(N_{Y/\mathbb{P}^M}(-H)) 
\end{equation}


does not surject. To see this note that (by sequence \ref{23} and diagram \ref{section}),  we have an exact sequence

$$0 \to \mathcal{O}_{\mathbb{P}^1}(-(b-e)) = p_*T_{Y/\mathbb{P}^1}(-H) \to p_*(T_{\mathbb{P}^M}|_Y(-H)) \to \operatorname{Ker}(\lambda) \to 0$$

Applying the functor $\operatorname{Hom}(p_*(-H-K_Y), _)$ we get 

\begin{multline}\label{30}
0 \to \operatorname{Hom}(p_*(-H-K_Y), 
 p_*(T_{\mathbb{P}^M}|_Y(-H)))  \to \operatorname{Hom}(p_*(-H-K_Y), p_*(T_{\mathbb{P}^M}|_Y(-H))) \\ \to \operatorname{Hom}(p_*(-H-K_Y),\operatorname{Ker}(\lambda)) \to \operatorname{Ext^1}(p_*(-H-K_Y), 
 p_*T_{Y/\mathbb{P}^1}(-H)) \\ \xrightarrow[]{\eta} \operatorname{Ext^1}(p_*(-H-K_Y), 
 p_*(T_{\mathbb{P}^M}|_Y(-H))
\end{multline}

%Since $p_*(-H-K_Y) = \mathcal{O}_{\mathbb{P}^1}(-(b-e-2))$, we have that $\operatorname{Hom}(\mathcal{O}_{\mathbb{P}^1}(-(b-e-2)), 
 %\mathcal{O}_{\mathbb{P}^1}(-(b-e))) = 0$ and $\operatorname{Ext^1}(\mathcal{O}_{\mathbb{P}^1}(-(b-e-2)), 
% \mathcal{O}_{\mathbb{P}^1}(-(b-e))) = $

We show that the map $\operatorname{Ext^1}(p_*(-H-K_Y), 
 p_*T_{Y/\mathbb{P}^1}(-H))  \xrightarrow[]{\eta} \operatorname{Ext^1}(p_*(-H-K_Y), 
 p_*(T_{\mathbb{P}^M}|_Y(-H))$ is an injective map and is in fact an isomorphism. To see this, push forward sequence \ref{24}, to get 

 \begin{equation}\label{25}
     0 \to \mathcal{O}_{\mathbb{P}^1}^{M+1} \to p_*(T_{\mathbb{P}^M}|_Y(-H)) \to R^1p_*(-H) \to 0
 \end{equation}

Note that $\operatorname{Ext^i}(p_*(-H-K_Y), \mathcal{O}_{\mathbb{P}^1}^{M+1}) = H^i(\mathcal{O}_{\mathbb{P}^1}(b-e-2))^{\oplus M+1} = 0$ for $i = 1,2$ if $b \geq e+1$. Hence applying the functor $\operatorname{Hom}(p_*(-H-K_Y), _)$ we get 

\begin{equation}\label{26}
     0 \to \operatorname{Ext^1}(p_*(-H-K_Y), 
 p_*(T_{\mathbb{P}^M}|_Y(-H))) \to \operatorname{Ext^1}(p_*(-H-K_Y), 
 R^1p_*(-H)) \to 0
 \end{equation}

Composing the map $\eta$ with the isomorphism in sequence \ref{25}, we see that it is enough to show that the following map (which we also call $\eta$ by a slight abuse of notation)

\begin{equation}\label{27}
   \operatorname{Ext^1}(p_*(-H-K_Y), 
 p_*T_{Y/\mathbb{P}^1}(-H))  \xrightarrow[]{\eta} \operatorname{Ext^1}(p_*(-H-K_Y), 
 R^1p_*(-H)) 
\end{equation}

is an isomorphism. To see this consider the relative Euler sequence corresponding to the projective bundle $Y \to \mathbb{P}^1$, twisted by $-H$ 

\begin{equation}\label{28}
0 \to \mathcal{O}_Y(-H) \to p^*(\mathcal{O}_{\mathbb{P}^1} \oplus \mathcal{O}_{\mathbb{P}^1}(-e))^* \otimes (C_0-H) \to T_{Y/\mathbb{P}^1}(-H) \to 0
\end{equation}

Pushing forward the above exact sequence, we have that 

\begin{equation}\label{29}
 0 \to p_*T_{Y/\mathbb{P}^1}(-H) \to R^1p_*(-H) \to 0 
\end{equation}

Now the map in sequence \ref{27} is induced by the isomorphism from sequence \ref{29} is hence an isomorphism. Going back to sequence \ref{30}, we have 

\begin{multline}\label{31}
0 \to \operatorname{Hom}(p_*(-H-K_Y), 
 p_*(T_{\mathbb{P}^M}|_Y(-H)))  \to \operatorname{Hom}(p_*(-H-K_Y), p_*(T_{\mathbb{P}^M}|_Y(-H))) \\ \to \operatorname{Hom}(p_*(-H-K_Y),\operatorname{Ker}(\lambda)) \to 0
\end{multline}

Therefore, to show that the map in \ref{32} does not surject, it is enough to show that

\begin{equation}\label{33}
  \operatorname{Hom}(p_*(-H-K_Y), p_*(T_{\mathbb{P}^M}|_Y(-H)) \to \operatorname{Hom}(p_*(-H-K_Y), p_*(N_{Y/\mathbb{P}^M}(-H))  
\end{equation}

is not surjective. This is the same as showing 

\begin{equation}\label{34}
   H^0(T_{\mathbb{P}^M}|_Y(K_Y))) \to H^0(N_{Y/\mathbb{P}^M}(K_Y)) 
\end{equation}
does not surject

Now by taking the cohomology of the exact sequence 

$$0 \to T_Y \otimes K_Y \to T_{\mathbb{P}^M|_Y} \otimes K_Y \to N_{Y/\mathbb{P}^M}(K_Y) \to 0$$


we have 

$$0 \to H^0(T_Y \otimes K_Y) \to H^0(T_{\mathbb{P}^M|_Y} \otimes K_Y) \to H^0(N_{Y/\mathbb{P}^M}(K_Y)) \to H^1(T_Y \otimes K_Y) \to H^1(T_{\mathbb{P}^M|_Y} \otimes K_Y)$$

Now the required non-surjection of sequence \ref{34} follows from the fact that $H^1(T_Y \otimes K_Y) = \mathbb{C}^2$ and $H^1(T_{\mathbb{P}^M|_Y} \otimes K_Y) = \mathbb{C}$. \QEDB




\begin{lemma}\label{2C_0+ef-H}
Let $H = aC_0+bf$ be a line bundle on $\mathbb{F}_e$ with $a \geq 1$. Then 
\begin{enumerate}
    \item[(1)] $H^1(2C_0+ef-H) = 0$ iff 
\[
\begin{cases}
   a = 1, & b \leq 1  \\
   a = 2, & b \leq e+1 \\
   a = 3 & \\
   a \geq 4, & b \geq ae-2e+1 \\
\end{cases}
\]

\item[(2)] $H^0(2C_0+ef-H) = 0$ iff 
\[
\begin{cases}
   a = 1, & b \geq e+1  \\
   a = 2, & b \geq e+1 \\
   a \geq 4 & \\
\end{cases}
\] 

\end{enumerate}

    
\end{lemma}

\begin{lemma}\label{2f-H}
Let $H = aC_0+bf$ be a line bundle on $\mathbb{F}_e$ with $a \geq 1$. Then 
\begin{enumerate}
    \item[(1)] $H^1(2f-H) = 0$ iff 
\[
\begin{cases}
   a = 1 &   \\
   a \geq 2, & b \geq ae-e+3 \\
   
\end{cases}
\]

\item[(2)] $H^0(2f-H) = 0$ 



\end{enumerate}  

\end{lemma}


\begin{lemma}\label{-H-K_Y}
Let $H = aC_0+bf$ be a line bundle on $\mathbb{F}_e$ with $a \geq 1$ and $b \geq ae+1$. Then $H^1(-H+K_{F_e}) = 0$. Further, $H^1(-H-K_{F_e}) = 0$ iff 
 \[
\begin{cases}
   a = 1, & b \leq 3  \\
   a = 2, & b \leq 3+e \\
   a = 3 & \\
   a \geq 4, & b \geq ae-2e+3 \\
\end{cases}
 \]

\end{lemma}



\begin{lemma}\label{-H-2K_Y}
Let $H = aC_0+bf$ be a line bundle on $\mathbb{F}_e$ with $a \geq 1$ and $b \geq ae+1$. Then $H^0(-H-2K_{\mathbb{F}_e}) = 0$ iff 
\[
\begin{cases}
   a \leq 4 & b \geq 2e+5  \\
   a \geq 5 &   \\
   
\end{cases}
\]
    
\end{lemma}

\begin{lemma}\label{T_Y(-H)}
Let $H = aC_0+bf$ be a line bundle on $\mathbb{F}_e$ with $a \geq 1$ and $b \geq ae+1$. Then 
\begin{enumerate}
    \item[(1)] $H^1(T_{\mathbb{F}_e}(-H)) = 0$ if 
\[
\begin{cases}
   a = 1, & b \leq 1  \\
   a = 2, & ae-e+3 \leq b \leq e+1  \\
   a = 3, & b \geq ae-e+3 \\
   a \geq 4, & b \geq \operatorname{max}\{ae-2e+3, ae-e+3\}
   
\end{cases}
\]

  \item[(2)] $H^1(T_{\mathbb{F}_e}(-H+K_{\mathbb{F}_e})) = 0$ if 
\[
\begin{cases}
   a = 1, & b \geq e+1  \\
   a \geq 2, & b \geq ae+1 \\
\end{cases}
\] 

\item[(3)] $H^0(T_{\mathbb{F}_e}(-H)) = 0$ if 
\[
\begin{cases}
   a = 1, & b \geq e+1  \\
   a = 2, & b \geq e+1 \\
   a \geq 3 & \\
\end{cases}
\] 

\end{enumerate}

\end{lemma}

\noindent\textit{Proof.} Let $Y = \mathbb{F}_e$ and let $p: Y \to \mathbb{P}^1$ denote the projection to $\mathbb{P}^1$. Note that we have an exact sequence 
$$0 \to T_{Y/\mathbb{P}^1} \to T_Y \to p^*T_{\mathbb{P}^1} \to 0$$
and that $T_{Y/\mathbb{P}^1} = 2C_0+ef$ and $p^*T_{\mathbb{P}^1} = 2f$. So it is enough to show $H^1(2C_0+ef-H) = H^1(2f-H) = 0$. Now the claim follows by combining Lemmas \ref{2C_0+ef-H} and \ref{2f-H}. \QEDA

\begin{proposition}\label{calculating beta}
%Let $V = Y_1 \bigcup_E Y_2 \hookrightarrow \mathbb{P}^M$ be a simple normal crossing variety of dimension $d > 1$, embedded by the complete linear series of a very ample line bundle $\widetilde{H}$, where $Y_1 = \mathbb{F}_e \hookrightarrow \mathbb{P}^N$, $e \leq 2$ is embedded by the complete linear series of a very ample line bundle $|aC_0+bf|$, $E \in |-K_{Y_1}|$, $Y_2$ is a general element of $\mathcal{F}_{E,Y}$ and $M = N + h^0((a-2)C_0+(b-e-2)f)$ (i.e, we are in the situation of Theorem \ref{Hirzebruch main}, $(1)$ (b). Then $\alpha(V) = \beta = 0$ if 
Let $\widetilde{Y} \hookrightarrow \mathbb{P}^M$ be a $K3$ carpet embedded by the complete linear series of a very ample line bundle $\widetilde{H}$ such that the embedding induced on the reduced part is $Y = \mathbb{F}_e \hookrightarrow \mathbb{P}^N$ embedded by the complete linear series of a very ample line bundle $H = aC_0+bf$.   
Then 

\begin{enumerate}
    \item $\beta = 0$ if
    \[
\begin{cases}
    a = 3, b \geq 5 & e = 0  \\
    a = 3, b \geq 7 & e = 1  \\
    a = 3, b \geq 9 & e = 2  \\
\color{blue}
    a = 3, b \geq 11 & e \geq 3 \\
    a = 3, b \geq 3e+1 & e \geq 4 \\
\color{black}
    a = 4, b \geq 5 & e = 0 \\
    a \geq 5, b \geq 3 & e = 0 \\
    a \geq 4, b \geq a+2 & e = 1 \\
    a \geq 4, b \geq ae+1 & e \geq 2 \\
    
\end{cases}
  \]

  \item If $a = 2$, then $\gamma = 0$ if $b \geq 2e+5$

  \item If $a = 1$, then $\alpha(\widetilde{Y}) = 0$ if $b \geq 2e+5$


\end{enumerate}
  

 


\end{proposition}
 
\noindent\textit{Proof $(1)$.} This follows by combining Lemma \ref{-H-K_Y}, \ref{-H-2K_Y}, \ref{T_Y(-H)} and applying Theorem \ref{ZL for K3 carpets}. \\

\noindent\textit{Proof $(2)$.} Note that by Lemma \ref{T_Y(-H)}, $(2)$, $h^1(T_Y(-H+K_Y)) = 0$ if $a = 2$ and $b \geq 2e+1$. Similarly by Lemma \ref{-H-2K_Y}, $h^0(-H-2K_Y) = 0$ if $a = 2$ and $b \geq 2e+5$. Now we show that $h^0(N_{Y/\mathbb{P}^M}(-H)) = M+1$. Twisting the sequence $$0 \to T_Y \to T_{\mathbb{P}^M}|_Y \to N_{Y/\mathbb{P}^M} \to 0$$ by $-H$ and taking cohomology, we have 

\begin{equation}\label{35}
0 \to H^0(T_Y(-H)) \to H^0(T_{\mathbb{P}^M}|_Y(-H)) \to H^0(N_{Y/\mathbb{P}^M}(-H)) \to H^1(T_Y(-H)) \to H^1(T_{\mathbb{P}^M}|_Y(-H))
\end{equation}

Now by Lemma \ref{T_Y(-H)}, part $(3)$, we have that $H^0(T_Y(-H)) = 0$ if $a = 2$ and $b \geq e+1$. Turning to the Euler sequence on $\mathbb{P}^M$, restricting to $Y$ twisting by $-H$ and taking cohomology, we get

\begin{equation}\label{36}
    0 \to H^0(\mathcal{O}_Y^{\oplus M+1}) \to H^0(T_{\mathbb{P}^M}|_Y(-H)) \to 0
\end{equation}

Hence $h^0(T_{\mathbb{P}^M}|_Y(-H)) = M+1$. So it is enough to show that the map $H^1(T_Y(-H)) \to H^1(T_{\mathbb{P}^M}|_Y(-H))$ is injective. Continuing with the long exact sequence of cohomology in \ref{36}, we have that 

\begin{equation}\label{36}
    0 \to H^1(T_{\mathbb{P}^M}|_Y(-H)) \to H^2(-H) \to 0
\end{equation}

Similary tensoring the exact sequence of horizontal and vertical tangent bundles by $-H$, we have

\begin{equation}\label{37}
    0 \to T_{Y/\mathbb{P}^1}(-H) \to T_Y(-H) \to p^*(T_{\mathbb{P}^1})(-H) \to 0
\end{equation}

and taking the cohomology, we have that 

\begin{equation}\label{38}
    H^0(p^*(T_{\mathbb{P}^1})(-H)) \to H^1(T_{Y/\mathbb{P}^1}(-H)) \to H^1(T_Y(-H)) \to H^1(p^*(T_{\mathbb{P}^1})(-H)) \to 0
\end{equation}

Now note that $p^*(T_{\mathbb{P}^1})(-H) = -2C_0-(b-2)f$. We have that $H^0(-2C_0-(b-2)f) = 0$ and $H^1(-2C_0-(b-2)f) = H^1((b-e-4)f)^* = 0$ if $b \geq e+3$. hence we have 

\begin{equation}\label{39}
    0 \to H^1(T_{Y/\mathbb{P}^1}(-H)) \to H^1(T_Y(-H)) \to  0
\end{equation}

Composing the map $H^1(T_Y(-H)) \to H^1(T_{\mathbb{P}^M}|_Y(-H))$ in sequence \ref{35} with the isomorphisms given by the sequences \ref{36} and \ref{39}, we see that it is enough to show that the map 

\begin{equation}\label{40}
     H^1(T_{Y/\mathbb{P}^1}(-H)) \to H^2(-H)
\end{equation}

is injective. But the above map comes as follows: Take the relative Euler sequence of the map $Y \to \mathbb{P}^1$, given in sequence \ref{28} twist by $-H$ and take the cohomology to get

\begin{equation}\label{41}
    0 \to H^1(T_{Y/\mathbb{P}^1}(-H)) \to H^2(-H)
\end{equation}

Hence we have that $\gamma = 0$ if $a = 2$ and $b \geq 2e+5$. \\

\noindent\textit{Proof $(3)$.} We finally settle the case $a = 1$.Tensoring sequence \ref{1} by $-\widetilde{H}$ and taking the cohomology, we have 

\begin{multline}\label{42}
0 \to H^0(N_{\widetilde{Y}/\mathbb{P}^M} (-H+K_Y)) \to H^0(N_{\widetilde{Y}/\mathbb{P}^M}(-\widetilde{H})) \to H^0(N_{\widetilde{Y}/\mathbb{P}^M}(-H))  \to H^1(N_{\widetilde{Y}/\mathbb{P}^M} (-H+K_Y)) \to H^1(N_{\widetilde{Y}/\mathbb{P}^M}(-\widetilde{H})) \\ \to H^1(N_{\widetilde{Y}/\mathbb{P}^M}(-H)) \to H^2(N_{\widetilde{Y}/\mathbb{P}^M} (-H+K_Y)) \to H^2(N_{\widetilde{Y}/\mathbb{P}^M}(-\widetilde{H})) \to H^2(N_{\widetilde{Y}/\mathbb{P}^M}(-H)) \to 0  
\end{multline}

We compute $H^i(N_{\widetilde{Y}/\mathbb{P}^M} (-H+K_Y)) = 0$ for $i = 0,1,2$. Tensoring sequence \ref{2} by $-H+K_Y$ and taking the cohomology, we have that 

\begin{multline}\label{43}
   0 \to H^0(\mathcal{H}om(I_{\widetilde{Y}}/I_{Y}^2, \mathcal{O}_Y)(-H+K_Y)) \to H^0(N_{\widetilde{Y}/\mathbb{P}^M}(-H+K_Y))  \to H^0(\mathcal{O}_Y(-H-K_Y)) \to H^1(\mathcal{H}om(I_{\widetilde{Y}}/I_{Y}^2, \mathcal{O}_Y)(-H+K_Y)) \\ \to H^1(N_{\widetilde{Y}/\mathbb{P}^M}(-H+K_Y))  \to H^1(\mathcal{O}_Y(-H-K_Y)) \to H^2(\mathcal{H}om(I_{\widetilde{Y}}/I_{Y}^2, \mathcal{O}_Y)(-H+K_Y)) \to H^2(N_{\widetilde{Y}/\mathbb{P}^M}(-H+K_Y)) \\  \to H^2(\mathcal{O}_Y(-H-K_Y)) \to 0
\end{multline}

Note that since $H = C_0+bf$, we have that $H^0(-H-K_Y) = 0$ if $b \geq e+3$. Further, $h^2(-H-K_Y) = h^0(H+2K_Y) = h^0(-3C_0+(b-2e-4)f) = 0$.

Tensoring sequence \ref{3} by $-H+K_Y$ and taking the cohomology, we have that 

\begin{multline}\label{44}
 0 \to H^0(-H) = (0) \to H^0(N_{Y/\mathbb{P}^M}(-H+K_Y)) \to H^0(\mathcal{H}om(I_{\widetilde{Y}}/I_{Y}^2, \mathcal{O}_Y)(-H+K_Y)) \to H^1(-H) = (0) \\ \to H^1(N_{Y/\mathbb{P}^M}(-H+K_Y)) \to H^1(\mathcal{H}om(I_{\widetilde{Y}}/I_{Y}^2, \mathcal{O}_Y)(-H+K_Y)) \to H^2(-H) = (0) \\ \to
 H^2(N_{Y/\mathbb{P}^M}(-H+K_Y)) \to H^2(\mathcal{H}om(I_{\widetilde{Y}}/I_{Y}^2, \mathcal{O}_Y)(-H+K_Y)) \to 0
\end{multline}

Now using the Euler exact sequence pulled back to $Y$  and tensored with $-H+K_Y$, we have that 

\begin{multline}\label{45}
 0 \to H^0(-H+K_Y) \to H^0(K_Y)^{\oplus M+1} = (0) \to H^0(T_{\mathbb{P}^M}|_Y (-H+K_Y)) \to H^1(-H+K_Y) \to H^1(K_Y)^{\oplus M+1} \\ \to H^1(T_{\mathbb{P}^M}|_Y (-H+K_Y)) \to H^2(-H+K_Y) \to H^2(K_Y)^{\oplus M+1} \to H^2(T_{\mathbb{P}^M}|_Y (-H+K_Y)) \to 0   
\end{multline}

We have $H^1(-H+K_Y) = 0$ for $b \geq e+1$ by Lemma \ref{-H-K_Y} and $H^1(K_Y) = 0$. Further the map $H^2(-H+K_Y) \to H^2(K_Y)^{\oplus M+1} = H^2(K_Y) \otimes H^0(H)$ is an isomorphism since it is the dual of the isomorphism $H^0(\mathcal{O}_Y) \otimes H^0(H) \to H^0(H)$. Putting all this together in sequence \ref{45}, we have

\begin{equation}\label{46}
  H^i(T_{\mathbb{P}^M}|_Y (-H+K_Y)) = 0, i = 0,1,2, b \geq e+1
\end{equation}

On the other hand, from the sequence 

\begin{equation}
    0 \to T_{Y/\mathbb{P}^1}(-H+K) = -C_0-(b+2)f \to T_Y(-H+K) \to p^*T_{\mathbb{P}^1}(-H+K) = -3C_0-(b+e)f \to 0
\end{equation}

we have that if $b \geq e+1$

\begin{equation}\label{47}
    H^{i}(T_Y(-H+K_Y)) = 0, i = 0,1, H^2(T_Y(-H+K_Y)) = H^0(C_0+(b-2))^*
\end{equation}

Hence from the sequence 

\begin{equation*}
    0 \to T_Y(-H+K_Y) \to T_{\mathbb{P}^M}|_Y(-H+K_Y)) \to N_{Y/\mathbb{P}^M}(-H+K_Y) \to 0
\end{equation*}

we conclude that for $b \geq e+1$, 

\begin{equation}\label{48}
    H^0(N_{Y/\mathbb{P}^M}(-H+K_Y)) = 0, H^1(N_{Y/\mathbb{P}^M}(-H+K_Y)) =  H^{2}(T_Y(-H+K_Y)) = H^0(C_0+(b-2)f)^*,  H^2(N_{Y/\mathbb{P}^M}(-H+K_Y)) = 0
\end{equation}

Plugging everything back into sequence \ref{44}, we have that 

\begin{multline}\label{49}
 H^0(\mathcal{H}om(I_{\widetilde{Y}}/I_{Y}^2, \mathcal{O}_Y)(-H+K_Y)) = 0, H^1(\mathcal{H}om(I_{\widetilde{Y}}/I_{Y}^2, \mathcal{O}_Y)(-H+K_Y)) = H^{2}(T_Y(-H+K_Y)) = H^0(C_0+(b-2))^*, \\  H^2(\mathcal{H}om(I_{\widetilde{Y}}/I_{Y}^2, \mathcal{O}_Y)(-H+K_Y)) = 0 
\end{multline}

Going back to sequence \ref{43}, and noting that $H^0(-H-K_Y) = 0$ for $b \geq e+3$, we have that for $b \geq e+3$, 
%for $b \geq 2e+5$

\begin{equation}\label{50}
   H^0(N_{\widetilde{Y}/\mathbb{P}^M}(-H+K_Y)) = 0, H^2(N_{\widetilde{Y}/\mathbb{P}^M}(-H+K_Y)) = 0
\end{equation}

and an exact sequence 

\begin{equation}\label{51}
0 \to H^2(T_Y(-H+K_Y)) \to  H^1(N_{\widetilde{Y}/\mathbb{P}^M}(-H+K_Y)) \to H^1(-H-K_Y) \to 0   
\end{equation}

yielding $h^1(N_{\widetilde{Y}/\mathbb{P}^M}(-H+K_Y)) = h^2(T_Y(-H+K_Y) + h^1(-H-K_Y)$

Now we compute $H^i(N_{\widetilde{Y}/\mathbb{P}^M} (-H)) = 0$ for $i = 0,1,2$. Tensoring sequence \ref{2} by $-H$ and taking the cohomology, we have that 

\begin{multline}\label{52}
   0 \to H^0(\mathcal{H}om(I_{\widetilde{Y}}/I_{Y}^2, \mathcal{O}_Y)(-H)) \to H^0(N_{\widetilde{Y}/\mathbb{P}^M}(-H))  \to H^0(\mathcal{O}_Y(-H-2K_Y)) \to H^1(\mathcal{H}om(I_{\widetilde{Y}}/I_{Y}^2, \mathcal{O}_Y)(-H)) \\ \to H^1(N_{\widetilde{Y}/\mathbb{P}^M}(-H))  \to H^1(\mathcal{O}_Y(-H-2K_Y)) \to H^2(\mathcal{H}om(I_{\widetilde{Y}}/I_{Y}^2, \mathcal{O}_Y)(-H)) \to H^2(N_{\widetilde{Y}/\mathbb{P}^M}(-H)) \\  \to H^2(\mathcal{O}_Y(-H-2K_Y)) \to 0
\end{multline}

Note that since $H = C_0+bf$, we have that $H^0(-H-2K_Y) = 0$ if $b \geq 2e+5$.

Tensoring sequence \ref{3} by $-H$ and taking the cohomology, we have that 

\begin{multline}\label{53}
 0 \to H^0(-H-K_Y) = (0) \to H^0(N_{Y/\mathbb{P}^M}(-H)) \to H^0(\mathcal{H}om(I_{\widetilde{Y}}/I_{Y}^2, \mathcal{O}_Y)(-H)) \to H^1(-H-K_Y)  \\ \to H^1(N_{Y/\mathbb{P}^M}(-H)) \to H^1(\mathcal{H}om(I_{\widetilde{Y}}/I_{Y}^2, \mathcal{O}_Y)(-H)) \to H^2(-H-K_Y) = (0) \\ \to
 H^2(N_{Y/\mathbb{P}^M}(-H)) \to H^2(\mathcal{H}om(I_{\widetilde{Y}}/I_{Y}^2, \mathcal{O}_Y)(-H)) \to 0
\end{multline}

Now using the Euler exact sequence pulled back to $Y$  and tensored with $-H$, we have that 

\begin{multline}\label{54}
 0 \to H^0(-H) \to H^0(\mathcal{O}_Y)^{\oplus M+1} \to H^0(T_{\mathbb{P}^M}|_Y (-H)) \to H^1(-H) = (0) \to H^1(\mathcal{O}_Y)^{\oplus M+1} = (0) \\ \to H^1(T_{\mathbb{P}^M}|_Y (-H)) \to H^2(-H) \to H^2(\mathcal{O}_Y)^{\oplus M+1} \to H^2(T_{\mathbb{P}^M}|_Y (-H)) \to 0   
\end{multline}

We have $H^0(-H) = 0$, $H^1(-H) = 0$ by Kodaira vanishing theorem, $H^2(-H) = 0$ by pushing forward to $\mathbb{P}^1$  since $a = 1$.  Putting all this together in sequence \ref{54}, we have

\begin{equation}\label{55}
  h^0(T_{\mathbb{P}^M}|_Y (-H)) = M+1, h^1(T_{\mathbb{P}^M}|_Y (-H)) = h^2(T_{\mathbb{P}^M}|_Y (-H)) = 0 
\end{equation}

On the other hand, from the sequence 

\begin{equation}
    0 \to T_{Y/\mathbb{P}^1}(-H) = C_0-(b-e)f \to T_Y(-H) \to p^*T_{\mathbb{P}^1}(-H) = -C_0-(b-2)f \to 0
\end{equation}

we have that if $b \geq e-1$

\begin{equation}\label{56}
    H^{i}(T_Y(-H)) = 0, i = 0,2, H^1(T_Y(-H)) = H^1(C_0-(b-e)f)
\end{equation}

Hence from the sequence 

\begin{equation*}
    0 \to T_Y(-H) \to T_{\mathbb{P}^M}|_Y(-H)) \to N_{Y/\mathbb{P}^M}(-H) \to 0
\end{equation*}

we conclude that for $b \geq e-1$, 

\begin{equation}\label{57}
    H^1(N_{Y/\mathbb{P}^M}(-H)) = H^2(N_{Y/\mathbb{P}^M}(-H)) = 0
\end{equation}

and we have an exact sequence

\begin{equation}\label{58}
    0 \to H^0(T_{\mathbb{P}^M}|_Y (-H)) = \mathbb{C}^{M+1} \to H^0(N_{Y/\mathbb{P}^M}(-H)) \to H^1(T_Y(-H)) \to 0
\end{equation}

Plugging everything back into sequence \ref{53}, we have that 

\begin{equation}\label{59}
  H^1(\mathcal{H}om(I_{\widetilde{Y}}/I_{Y}^2, \mathcal{O}_Y)(-H)) =H^2(\mathcal{H}om(I_{\widetilde{Y}}/I_{Y}^2, \mathcal{O}_Y)(-H)) = 0 
\end{equation}

and an exact sequence

\begin{equation}\label{60}
    0 \to H^0(N_{Y/\mathbb{P}^M}(-H)) \to H^0(\mathcal{H}om(I_{\widetilde{Y}}/I_{Y}^2, \mathcal{O}_Y)(-H)) \to H^1(-H-K_Y) \to 0
\end{equation}

Going back to sequence \ref{52}, and noting that $H^0(-H-2K_Y) = 0$ for $b \geq 2e+5$, we have that for $b \geq 2e+5$

\begin{equation}\label{61}
   H^0(N_{\widetilde{Y}/\mathbb{P}^M}(-H)) = H^0(\mathcal{H}om(I_{\widetilde{Y}}/I_{Y}^2, \mathcal{O}_Y)(-H)), H^1(N_{\widetilde{Y}/\mathbb{P}^M}(-H)) = H^1(-H-2K_Y), H^2(N_{\widetilde{Y}/\mathbb{P}^M}(-H)) = 0
\end{equation}

Plugging \ref{50}, \ref{51}, \ref{61} into the exact sequence \ref{42}, we have 

\begin{align}\label{62}
   h^0(N_{\widetilde{Y}/\mathbb{P}^M}(-\widetilde{H}))  & = h^0(N_{\widetilde{Y}/\mathbb{P}^M}(-H)) - h^1(N_{\widetilde{Y}/\mathbb{P}^M}(-H+K_Y)) + h^1(N_{\widetilde{Y}/\mathbb{P}^M}(-\widetilde{H}))-h^1(N_{\widetilde{Y}/\mathbb{P}^M}(-H)) \\
    & = M+1 + h^1(T_Y(-H)) - h^2(T_Y(-H+K_Y)) + h^1(N_{\widetilde{Y}/\mathbb{P}^M}(-\widetilde{H}))-h^1(-H-2K_Y) 
\end{align}

Now note that by \ref{56}, $h^1(T_Y(-H)) = h^1(C_0-(b-e)f) = 2b-e-2$ and $h^2(T_Y(-H+K_Y)) = h^0(C_0+(b-2)f) = 2b-e-2$. Hence 

\begin{equation}\label{63}
   h^0(N_{\widetilde{Y}/\mathbb{P}^M}(-\widetilde{H})) = M+1 + h^1(N_{\widetilde{Y}/\mathbb{P}^M}(-\widetilde{H}))-h^1(-H-2K_Y) 
\end{equation}

We end the proof by showing $h^1(N_{\widetilde{Y}/\mathbb{P}^M}(-\widetilde{H}))=h^1(-H-2K_Y)$ if $b \geq e+3$ (\textcolor{blue}{The fact that we never use $b \geq 2e+5$ is not used from now is used crucially to get the correct lower genus bounds in the next section !}). Note that from sequence \ref{42} and equations \ref{50} and \ref{61}, we already have

\begin{multline}\label{64}
0 \to H^0(N_{\widetilde{Y}/\mathbb{P}^M}(-\widetilde{H})) \to H^0(N_{\widetilde{Y}/\mathbb{P}^M}(-H))  \xrightarrow[]{} H^1(N_{\widetilde{Y}/\mathbb{P}^M} (-H+K_Y)) \to H^1(N_{\widetilde{Y}/\mathbb{P}^M}(-\widetilde{H})) 
 \\ \to H^1(-H-2K_Y) \to 0  
\end{multline}

We show that the map $H^0(N_{\widetilde{Y}/\mathbb{P}^M}(-H))  \xrightarrow[]{} H^1(N_{\widetilde{Y}/\mathbb{P}^M} (-H+K_Y))$ is surjective. First note that we have the sequence 
$$0 \to H^0(\mathcal{H}om(I_{\widetilde{Y}}/I_{Y}^2, \mathcal{O}_Y)(-H)) \to H^0(N_{\widetilde{Y}/\mathbb{P}^M}(-H)) \to H^0(-H-2K_Y) \to 0$$

It is enough to show that the composed map $H^0(\mathcal{H}om(I_{\widetilde{Y}}/I_{Y}^2, \mathcal{O}_Y)(-H))  \xrightarrow[]{f} H^1(N_{\widetilde{Y}/\mathbb{P}^M} (-H+K_Y))$ is surjective
Note that from 
\ref{60}, we have 

$$0 \to H^0(N_{Y/\mathbb{P}^M}(-H)) \to  H^0(\mathcal{H}om(I_{\widetilde{Y}}/I_{Y}^2, \mathcal{O}_Y)(-H)) \xrightarrow[]{g} H^1(-H-K_Y) \to 0$$

while from \ref{48} and \ref{51}, we have the exact sequence 

$$ 0 \to H^1(N_{Y/\mathbb{P}^M}(-H+K_Y)) = H^2(T_Y(-H+K_Y)) \to  H^1(N_{\widetilde{Y}/\mathbb{P}^M}(-H+K_Y)) \xrightarrow[]{h} H^1(-H-K_Y) \to 0  $$

It is easy to see that $h \circ f = 2g$. 
%Since both maps $h \circ f$ and $g$ naturally arise out the data of the ribbon $\widetilde{Y} \subset \mathbb{P}^M$, or in other words chasing through the sequences \ref{1}, \ref{2} and \ref{3}, we have that $h \circ f = g$.
Therefore we have a commutative diagram of exact sequences

\[
\begin{tikzcd}\label{65}
 0 \arrow[r] & H^0(N_{Y/\mathbb{P}^M}(-H)) \arrow[r] \arrow[d, "a"] & H^0(N_{\widetilde{Y}/\mathbb{P}^M}(-H)) \arrow[r, "g"] \arrow[d, "f"] & H^1(-H-K_Y) \arrow[r] \arrow[d, "2\operatorname{id}"] & 0  \\
 0 \arrow[r] & H^1(N_{Y/\mathbb{P}^M}(-H+K_Y))  \arrow[r] & H^1(N_{\widetilde{Y}/\mathbb{P}^M}(-H+K_Y)) \arrow[r, "h"] & H^1(-H-K_Y) \arrow[r] & 0
\end{tikzcd}
\]

where $a$ is the map induced by $f$ on the kernel.
\color{blue}
%----------------------------------------------------------------------
% Compatibility of the coboundaries f, g, h.
% Replaces: "It is easy to see that h \circ f = g."
% Replace (3.1), (3.2), (3.3) below by \eqref{...} of your labels.
%----------------------------------------------------------------------

Throughout, $L := I_Y/I_{\widetilde Y}\cong K_Y$ denotes the ideal sheaf of $Y$ in
$\widetilde Y$, so that $0\to L\to\mathcal{O}_{\widetilde Y}\to\mathcal{O}_Y\to 0$ is exact
and $L\cdot L = 0$ in $\mathcal{O}_{\widetilde Y}$ (equivalently $I_Y\cdot L=0$).
Recall that $N_{\widetilde Y/\mathbb{P}^M}
=\mathcal{H}om_{\mathcal{O}_{\widetilde Y}}(I_{\widetilde Y}/I_{\widetilde Y}^2,\mathcal{O}_{\widetilde Y})$
is locally free on $\widetilde Y$, that
\[
N_{\widetilde Y/\mathbb{P}^M}\otimes\mathcal{O}_Y
  =\mathcal{H}om_{\mathcal{O}_Y}\big(I_{\widetilde Y}/I_Y I_{\widetilde Y},\mathcal{O}_Y\big),
\qquad
N_{\widetilde Y/\mathbb{P}^M}\otimes L
  =\mathcal{H}om_{\mathcal{O}_Y}\big(I_{\widetilde Y}/I_Y I_{\widetilde Y},L\big),
\]
and that (3.2) and (3.3) are the $\mathcal{O}_Y$-duals of
\[
0\to I_Y^2/I_Y I_{\widetilde Y}\to I_{\widetilde Y}/I_Y I_{\widetilde Y}\to I_{\widetilde Y}/I_Y^2\to 0,
\qquad
0\to I_{\widetilde Y}/I_Y^2\to I_Y/I_Y^2\to L\to 0,
\]
where $I_Y^2/I_Y I_{\widetilde Y}\cong L^{\otimes 2}$ via $\bar x\otimes\bar y\mapsto xy$;
this is the identification under which the quotient in (3.2) is
$\mathcal{H}om(L^{\otimes 2},\mathcal{O}_Y)=\mathcal{O}_Y(-2K_Y)$.

\begin{lemma}\label{lem:hfg}
Let
\begin{align*}
f&\colon H^0\big(\mathcal{H}om(I_{\widetilde Y}/I_Y^2,\mathcal{O}_Y)(-H)\big)
   \longrightarrow H^0\big(N_{\widetilde Y/\mathbb{P}^M}\otimes\mathcal{O}_Y(-H)\big)
   \xrightarrow{\ \partial\ } H^1\big(N_{\widetilde Y/\mathbb{P}^M}\otimes L(-H)\big),\\
h&\colon H^1\big(N_{\widetilde Y/\mathbb{P}^M}\otimes L(-H)\big)
   \longrightarrow H^1\big(\mathcal{H}om(L^{\otimes 2},L(-H))\big)=H^1\big(L^{-1}(-H)\big)=H^1(-H-K_Y),\\
g&\colon H^0\big(\mathcal{H}om(I_{\widetilde Y}/I_Y^2,\mathcal{O}_Y)(-H)\big)
   \xrightarrow{\ \partial\ } H^1\big(\mathcal{H}om(L,\mathcal{O}_Y(-H))\big)=H^1\big(L^{-1}(-H)\big)=H^1(-H-K_Y),
\end{align*}
where the first $\partial$ is the connecting homomorphism of
(3.1)$\,\otimes\,\mathcal{O}_{\widetilde Y}(-\widetilde H)$, the second is that of
(3.3)$\,\otimes\,\mathcal{O}_Y(-H)$, the first arrow in $f$ is induced by the injection in (3.2),
$h$ is induced by the surjection in (3.2)$\,\otimes\,L(-H)$, and the two target groups are
identified by $\theta\mapsto\theta\otimes\mathrm{id}_L$. Then
\[
h\circ f \;=\; 2\,g .
\]
\end{lemma}

\begin{proof}
We compute with \v{C}ech cocycles. Let $\{U_i\}$ be an affine open cover of $\widetilde Y$
(equivalently of $Y$: all sheaves involved are quasi-coherent and supported on $\widetilde Y$),
and put $U_{ij}=U_i\cap U_j$. Because the cover is affine, a section of the right-hand term of a
short exact sequence of quasi-coherent sheaves lifts over each $U_i$, and the connecting
homomorphism on $H^0$ is computed by lifting over $U_i$ and taking differences over $U_{ij}$.

Fix $\varphi\in H^0\big(\mathcal{H}om(I_{\widetilde Y}/I_Y^2,\mathcal{O}_Y)(-H)\big)$, i.e.\ an
$\mathcal{O}_Y$-linear map $\varphi\colon I_{\widetilde Y}/I_Y^2\to\mathcal{O}_Y(-H)$. Its image in
$H^0(N_{\widetilde Y/\mathbb{P}^M}\otimes\mathcal{O}_Y(-H))$ is $\bar\varphi:=\varphi\circ\pi$, where
$\pi\colon I_{\widetilde Y}/I_Y I_{\widetilde Y}\to I_{\widetilde Y}/I_Y^2$ is the projection. In particular
\begin{equation}\label{eq:kill}
\bar\varphi\,(I_Y^2)=0 .
\end{equation}

\emph{The cocycle for $f(\varphi)$.}
Since $N_{\widetilde Y/\mathbb{P}^M}$ is locally free, (3.1)$\,\otimes\,\mathcal{O}_{\widetilde Y}(-\widetilde H)$
is obtained by applying $\mathcal{H}om_{\mathcal{O}_{\widetilde Y}}(I_{\widetilde Y}/I_{\widetilde Y}^2,-)$ to
$0\to L(-H)\to\mathcal{O}_{\widetilde Y}(-\widetilde H)\to\mathcal{O}_Y(-H)\to 0$. Choose
$\mathcal{O}_{\widetilde Y}$-linear lifts
\[
\Phi_i\colon I_{\widetilde Y}/I_{\widetilde Y}^2\big|_{U_i}\longrightarrow
\mathcal{O}_{\widetilde Y}(-\widetilde H)\big|_{U_i}
\qquad\text{of}\qquad \bar\varphi|_{U_i}.
\]
Then $\Phi_{ij}:=\Phi_i-\Phi_j$ takes values in $L(-H)$, hence is a section over $U_{ij}$ of
$\mathcal{H}om_{\mathcal{O}_Y}(I_{\widetilde Y}/I_Y I_{\widetilde Y},L(-H))=N_{\widetilde Y/\mathbb{P}^M}\otimes L(-H)$,
and $f(\varphi)=[\{\Phi_{ij}\}]$. The map $h$ is restriction along
$L^{\otimes 2}\cong I_Y^2/I_Y I_{\widetilde Y}\hookrightarrow I_{\widetilde Y}/I_Y I_{\widetilde Y}$, so
\begin{equation}\label{eq:hf}
h(f(\varphi))=\big[\{\Phi_{ij}|_{L^{\otimes 2}}\}\big],
\qquad
\Phi_{ij}|_{L^{\otimes 2}}(\bar x\otimes\bar y)=\Phi_i(xy)-\Phi_j(xy)
\end{equation}
for local sections $x,y$ of $I_Y$.

\emph{The cocycle for $g(\varphi)$.}
Choose $\mathcal{O}_Y$-linear lifts $\psi_i\colon I_Y/I_Y^2|_{U_i}\to\mathcal{O}_Y(-H)|_{U_i}$ of
$\varphi|_{U_i}$ along the surjection in (3.3)$\,\otimes\,\mathcal{O}_Y(-H)$, i.e.\
$\psi_i|_{I_{\widetilde Y}/I_Y^2}=\varphi$. Then $\psi_{ij}:=\psi_i-\psi_j$ vanishes on
$I_{\widetilde Y}/I_Y^2$, hence factors through $L=I_Y/I_{\widetilde Y}$, and
$g(\varphi)=[\{\psi_{ij}\}]\in H^1(\mathcal{H}om(L,\mathcal{O}_Y(-H)))$.

\emph{The identification of the targets.}
Since $L$ is a line bundle, for $\theta\in\mathcal{H}om(L,\mathcal{O}_Y(-H))$ and local sections
$u=\alpha s$, $v=\beta s$ of $L$ ($s$ a local generator) one has
\begin{equation}\label{eq:sym}
\theta(u)\otimes v=\alpha\beta\,\theta(s)\otimes s=\theta(v)\otimes u
\qquad\text{in } \mathcal{O}_Y(-H)\otimes L=L(-H),
\end{equation}
so $\theta\otimes\mathrm{id}_L\in\mathcal{H}om(L^{\otimes 2},L(-H))$ is symmetric and does not depend
on the factor to which $\theta$ is applied.

\emph{Comparison.}
For each $i$ define
\[
B_i\colon I_Y\times I_Y\longrightarrow L(-H)\big|_{U_i},\qquad B_i(x,y):=\Phi_i(xy).
\]
This makes sense: $xy\in I_Y^2\subset I_{\widetilde Y}$, and by \eqref{eq:kill}
$\Phi_i(I_Y^2)\subset\ker\big(\mathcal{O}_{\widetilde Y}(-\widetilde H)\to\mathcal{O}_Y(-H)\big)=L(-H)$.
We claim:
\begin{enumerate}
\item[(i)] $B_i$ is symmetric and $\mathcal{O}_{\mathbb{P}^M}$-bilinear, and $B_i(x,y)=0$ whenever
$x\in I_Y^2$ or $y\in I_Y^2$. Hence $B_i$ descends to a symmetric $\mathcal{O}_Y$-bilinear form on
$I_Y/I_Y^2$ with values in $L(-H)$.
\item[(ii)] If $x\in I_{\widetilde Y}$, then $B_i(x,y)=\varphi(\bar x)\otimes\bar y$, where
$\bar x\in I_{\widetilde Y}/I_Y^2$, $\bar y\in L=I_Y/I_{\widetilde Y}$ and the product is taken in
$\mathcal{O}_Y(-H)\otimes L=L(-H)$. Symmetrically, if $y\in I_{\widetilde Y}$ then
$B_i(x,y)=\varphi(\bar y)\otimes\bar x$. In particular these values do not depend on $i$.
\end{enumerate}
For (i): symmetry is $xy=yx$, and bilinearity is the $\mathcal{O}_{\widetilde Y}$-linearity of $\Phi_i$
(a section $a$ of $\mathcal{O}_{\mathbb{P}^M}$ acts on $L(-H)$ through its image in $\mathcal{O}_Y$).
If $x=\sum_k x'_k x''_k$ with $x'_k,x''_k\in I_Y$, then $x'_k y\in I_Y^2\subset I_{\widetilde Y}$, so by
$\mathcal{O}_{\widetilde Y}$-linearity $\Phi_i(x'_k x''_k y)=\bar x''_k\cdot\Phi_i(x'_k y)$ with
$\bar x''_k\in L\subset\mathcal{O}_{\widetilde Y}$ and $\Phi_i(x'_k y)\in L(-H)$; since
$L\cdot L(-H)=0$ inside $\mathcal{O}_{\widetilde Y}(-\widetilde H)$, this vanishes.
For (ii): if $x\in I_{\widetilde Y}$, then $xy=y\cdot x$ with $x$ a section of $I_{\widetilde Y}$, so
$\Phi_i(xy)=\bar y\cdot\Phi_i(x)$ with $\bar y\in L\subset\mathcal{O}_{\widetilde Y}$. Multiplication by
$\bar y$ on $\mathcal{O}_{\widetilde Y}(-\widetilde H)$ kills $L(-H)$ and has image in $L(-H)$, hence
factors as $\mathcal{O}_{\widetilde Y}(-\widetilde H)\to\mathcal{O}_Y(-H)\xrightarrow{\;\cdot\,\bar y\;}L(-H)$.
As $\Phi_i(x)\equiv\bar\varphi(x)=\varphi(\bar x)$ modulo $L(-H)$, we get
$B_i(x,y)=\varphi(\bar x)\otimes\bar y$.

Next, using the lifts $\psi_i$, define the symmetric $\mathcal{O}_Y$-bilinear form
\[
B'_i\colon I_Y/I_Y^2\times I_Y/I_Y^2\longrightarrow L(-H)\big|_{U_i},\qquad
B'_i(x,y):=\psi_i(x)\otimes\bar y+\psi_i(y)\otimes\bar x .
\]
If $x\in I_{\widetilde Y}/I_Y^2$, then $\bar x=0$ in $L$ and $\psi_i(x)=\varphi(x)$, so
$B'_i(x,y)=\varphi(x)\otimes\bar y=B_i(x,y)$ by (ii); symmetrically for $y\in I_{\widetilde Y}/I_Y^2$.
Therefore
\[
C_i:=B_i-B'_i
\]
is a symmetric bilinear form on $I_Y/I_Y^2$ which vanishes as soon as one of its arguments lies
in $I_{\widetilde Y}/I_Y^2$. It therefore descends to a bilinear form on $L\times L$, i.e.\ to a section
\[
C_i\in\Gamma\big(U_i,\mathcal{H}om(L^{\otimes 2},L(-H))\big)=\Gamma\big(U_i,L^{-1}(-H)\big),
\]
a \v{C}ech $0$-cochain of $L^{-1}(-H)$.

Now work over $U_{ij}$ and regard everything as forms on $L\times L$. By \eqref{eq:hf} and (ii),
$\Phi_{ij}|_{L^{\otimes 2}}$ is the form $(\bar x,\bar y)\mapsto B_i(x,y)-B_j(x,y)$, which represents
$h(f(\varphi))$. On the other hand
\[
B_i-B_j=(C_i-C_j)+(B'_i-B'_j),
\qquad
(B'_i-B'_j)(u,v)=\psi_{ij}(u)\otimes v+\psi_{ij}(v)\otimes u=2\,\psi_{ij}(u)\otimes v
\]
for local sections $u,v$ of $L$, the last equality by \eqref{eq:sym}. Thus, as \v{C}ech $1$-cocycles
of $L^{-1}(-H)$,
\[
\{B_i-B_j\}=\delta\{C_i\}+2\,\{\psi_{ij}\otimes\mathrm{id}_L\},
\]
and passing to cohomology gives $h(f(\varphi))=2\,g(\varphi)$.
\end{proof}

\begin{remark}
Lemma~\ref{lem:hfg} uses nothing about $a,b,e$ or about $\widetilde Y$ being a K3 carpet: it holds
for any ribbon $\widetilde Y\subset\mathbb{P}^M$ on a smooth $Y$ with conormal bundle $L$ and locally
free $N_{\widetilde Y/\mathbb{P}^M}$. Since we work in characteristic zero, $2$ is a unit, so
$\operatorname{Im}(h\circ f)=\operatorname{Im}(g)$. Consequently the diagram following (3.65) commutes
once its right-hand vertical arrow is taken to be $2\cdot\mathrm{id}$ instead of $\mathrm{id}$. As this
arrow is still an isomorphism, the four lemma (or a direct chase) gives
$\operatorname{coker}(a)\cong\operatorname{coker}(f)$, so $f$ is surjective if and only if $a$ is.
\end{remark}
\color{black}
Therefore to show that $f$ is surjective, it is enough to show that $a$ is surjective. Now from \ref{58}, we have an exact sequence
\begin{equation}\label{66}
0 \to H^0(T_{\mathbb{P}^M}|_Y (-H)) = \mathbb{C}^{M+1} \to H^0(N_{Y/\mathbb{P}^M}(-H)) \to H^1(T_Y(-H)) \to 0
\end{equation}
while from \ref{48}, we have

\begin{equation}\label{67}
0 \to H^1(T_{\mathbb{P}^M}|_Y (-H+K_Y)) = (0) \to H^1(N_{Y/\mathbb{P}^M}(-H+K_Y)) \to H^2(T_Y(-H+K_Y) \to 0
\end{equation}

Now we claim that $$H^1(T_Y(-H)) = H^2(T_Y(-H+K_Y) = H^1(\mathcal{O}_{\mathbb{P}^1}(-b) \oplus \mathcal{O}_{\mathbb{P}^1}(-(b-e)))$$

First note the exact sequence 

$$0 \to T_{Y/\mathbb{P}^1}(-H) = C_0-(b-e)f \to T_Y(-H) \to p^*T_{\mathbb{P}^1}(-H) = -C_0-(b-2)f \to 0 $$

when pushed forward along the map $Y \to \mathbb{P}^1$ and noting that $p_*(C_0-(b-e)f) = \mathcal{O}_{\mathbb{P}^1}(-b) \oplus \mathcal{O}_{\mathbb{P}^1}(-(b-e))$, $R^1p_*(C_0-(b-e)f) = 0$, $p_*(-C_0-(b-2)f) = R^1p_*(-C_0-(b-2)f) = 0$ we have 

\begin{equation}\label{68}
 p_*T_Y(-H) = \mathcal{O}_{\mathbb{P}^1}(-b) \oplus \mathcal{O}_{\mathbb{P}^1}(-(b-e)), R^1p_*T_Y(-H) = 0 
\end{equation}

Hence $H^1(T_Y(-H)) = H^1(p_*T_Y(-H)) = H^1(\mathcal{O}_{\mathbb{P}^1}(-b) \oplus \mathcal{O}_{\mathbb{P}^1}(-(b-e)))$

Next note the exact sequence 

$$0 \to T_{Y/\mathbb{P}^1}(-H+K_Y) = -C_0-(b+2)f \to T_Y(-H) \to p^*T_{\mathbb{P}^1}(-H) = -3C_0-(b+e)f \to 0 $$

when pushed forward along the map $Y \to \mathbb{P}^1$ and noting that $p_*(-C_0-(b+2)f) = R^1p_*(-C_0-(b+2)f) = 0$, $p_*(-3C_0-(b+e)f) = 0$ and $R^1p_*(-3C_0-(b+e)f) = \mathcal{O}_{\mathbb{P}^1}(-b) \oplus \mathcal{O}_{\mathbb{P}^1}(-(b-e))$ we have 

\begin{equation}\label{69}
 p_*T_Y(-H+K_Y) = 0, R^1p_*T_Y(-H+K_Y) =  \mathcal{O}_{\mathbb{P}^1}(-b) \oplus \mathcal{O}_{\mathbb{P}^1}(-(b-e))
\end{equation}

Hence $H^2(T_Y(-H+K_Y)) = H^1(R^1p_*T_Y(-H+K_Y)) = H^1(\mathcal{O}_{\mathbb{P}^1}(-b) \oplus \mathcal{O}_{\mathbb{P}^1}(-(b-e)))$

Using \ref{66}, \ref{67}, \ref{68}, \ref{69} we have a diagram of exact sequences

\[
\begin{tikzcd}
  0 \arrow[r] & H^0(T_{\mathbb{P}^M}|_Y (-H)) = \mathbb{C}^{N+1} \arrow[r] \arrow[d, "d"] & H^0(N_{Y/\mathbb{P}^M}(-H)) \arrow[r, "c"] \arrow[d, "a"] & H^1(\mathcal{O}_{\mathbb{P}^1}(-b) \oplus \mathcal{O}_{\mathbb{P}^1}(-(b-e))) \arrow[r] & 0 \\
  0 \arrow[r] & H^1(T_{\mathbb{P}^M}|_Y (-H+K_Y)) = (0) \arrow[r]  & H^1(N_{Y/\mathbb{P}^M}(-H+K_Y)) \arrow[r, "b"]  & H^1(\mathcal{O}_{\mathbb{P}^1}(-b) \oplus \mathcal{O}_{\mathbb{P}^1}(-(b-e))) \arrow[r] & 0
\end{tikzcd}
\]

The map $d$ is obtained as follows:
Tensor the exact sequence 
$$0 \to K_Y \to \mathcal{O}_{\widetilde{Y}} \to \mathcal{O}_Y \to 0$$ with $T_{\mathbb{P}^M}|_{\widetilde{Y}}(-\widetilde{H})$ and taking cohomology to get 

$$0 \to H^0(T_{\mathbb{P}^M}|_Y (-H+K_Y)) \to H^0(T_{\mathbb{P}^M}|_{\widetilde{Y}}(-\widetilde{H})) \to H^0(T_{\mathbb{P}^M}|_Y (-H)) \xrightarrow[]{d} H^1(T_{\mathbb{P}^M}|_Y (-H+K_Y))  $$

This in turn induces the map $e$ in the diagram 

\[
\begin{tikzcd}\label{70}
  0 \arrow[r] & H^0(T_{\mathbb{P}^M}|_Y (-H)) = \mathbb{C}^{N+1} \arrow[r] \arrow[d, "d"]  & H^0(N_{Y/\mathbb{P}^M}(-H)) \arrow[r, "c"] \arrow[d, "a"] & H^1(\mathcal{O}_{\mathbb{P}^1}(-b) \oplus \mathcal{O}_{\mathbb{P}^1}(-(b-e))) \arrow[r] \arrow[d, "e"] & 0 \\
  0 \arrow[r] & H^1(T_{\mathbb{P}^M}|_Y (-H+K_Y)) = (0) \arrow[r]  & H^1(N_{Y/\mathbb{P}^M}(-H+K_Y)) \arrow[r, "b"]  & H^1(\mathcal{O}_{\mathbb{P}^1}(-b) \oplus \mathcal{O}_{\mathbb{P}^1}(-(b-e))) \arrow[r] & 0
\end{tikzcd}
\]

\color{blue}
We describe the map $e$ and show that it is an isomorphism. By the projection formula of Ext groups (see \cite[\href{https://stacks.math.columbia.edu/tag/08JM}{Tag 08JM}]{stacks-project}), we have 
$$H^1(T_Y(-H)) = \operatorname{Ext}^1(\Omega_Y(H),\mathcal{O}_Y) = \operatorname{Ext}^1((Rp_*T_Y(-H))^*,\mathcal{O}_{\mathbb{P}^1})$$

By the spectral sequence and by \Cref{68}, the term $\operatorname{Ext}^1((Rp_*T_Y(-H))^*,\mathcal{O}_{\mathbb{P}^1})$ has contributions from 
$$\operatorname{Ext}^1((p_*T_Y(-H))^*,\mathcal{O}_{\mathbb{P}^1}) =  \operatorname{Ext}^1(\mathcal{O}_{\mathbb{P}^1}(b) \oplus \mathcal{O}_{\mathbb{P}^1}(b-e) ,\mathcal{O}_{\mathbb{P}^1}) \operatorname{and} \operatorname{Hom}((R^1p_*T_Y(-H))^*,\mathcal{O}_{\mathbb{P}^1}) = 0$$

Similarly, 

$$H^2(T_Y(-H+K_Y)) = \operatorname{Ext}^2(\Omega_Y(H-K_Y),\mathcal{O}_Y) = \operatorname{Ext}^2((Rp_*T_Y(-H+K))^*,\mathcal{O}_{\mathbb{P}^1})$$

and by the spectral sequence and by \Cref{69} the term $\operatorname{Ext}^2((Rp_*T_Y(-H+K_Y))^*,\mathcal{O}_{\mathbb{P}^1})$ has contributions from 
$$\operatorname{Ext}^2((p_*T_Y(-H+K_Y))^*,\mathcal{O}_{\mathbb{P}^1}) = 0, \operatorname{and} \operatorname{Ext}^1((R^1p_*T_Y(-H))^*, \mathcal{O}_{\mathbb{P}^1}) = \operatorname{Ext}^1(\mathcal{O}_{\mathbb{P}^1}(b) \oplus \mathcal{O}_{\mathbb{P}^1}(b-e) ,\mathcal{O}_{\mathbb{P}^1})$$


On the other hand the space of ribbons is given by

$$H^1(T_Y(K_Y)) = \operatorname{Ext}^1(\Omega_Y(-K_Y),\mathcal{O}_Y) = \operatorname{Ext}^1((Rp_*T_Y(K))^*,\mathcal{O}_{\mathbb{P}^1})$$
which once again by the spectral sequence is 
$$\operatorname{Ext}^1((p_*T_Y(K))^*,\mathcal{O}_{\mathbb{P}^1}) \oplus \operatorname{Hom}((R^1p_*T_Y(K))^*,\mathcal{O}_{\mathbb{P}^1})$$
By \cite[Proposition $1.7$, page $2481$]{GP97}, the above is 
$$\operatorname{Ext}^1(\mathcal{O}_{\mathbb{P}^1}(2),\mathcal{O}_{\mathbb{P}^1}) \oplus \operatorname{Hom}(\mathcal{O}_{\mathbb{P}^1},\mathcal{O}_{\mathbb{P}^1})$$
Now we have a Yoneda pairing of Ext groups

\begin{equation}\label{yoneda}
(\operatorname{Ext}^1(\mathcal{O}_{\mathbb{P}^1}(2), \mathcal{O}_{\mathbb{P}^1}) \oplus  \operatorname{Hom}(\mathcal{O}_{\mathbb{P}^1},\mathcal{O}_{\mathbb{P}^1})) \times \operatorname{Ext}^1(\mathcal{O}_{\mathbb{P}^1}(b) \oplus \mathcal{O}_{\mathbb{P}^1}(b-e),\mathcal{O}_{\mathbb{P}^1}) \to \operatorname{Ext}^1(\mathcal{O}_{\mathbb{P}^1}(b) \oplus \mathcal{O}_{\mathbb{P}^1}(b-e),\mathcal{O}_{\mathbb{P}^1}) 
\end{equation}

For a non-split ribbon $\widetilde{Y} \in \operatorname{Ext}^1((Rp_*T_Y(K))^*,\mathcal{O}_{\mathbb{P}^1})$, the image if its class in $\operatorname{Hom}(\mathcal{O}_{\mathbb{P}^1},\mathcal{O}_{\mathbb{P}^1})$ is non-zero and the map $e$ is given by the map induced by \Cref{yoneda} when we fix the image of the class of the ribbon in $\operatorname{Hom}(\mathcal{O}_{\mathbb{P}^1},\mathcal{O}_{\mathbb{P}^1})$
$$\operatorname{Ext}^1(\mathcal{O}_{\mathbb{P}^1}(-b) \oplus \mathcal{O}_{\mathbb{P}^1}(-(b-e)),\mathcal{O}_{\mathbb{P}^1}) \to \operatorname{Ext}^1(\mathcal{O}_{\mathbb{P}^1}(-b) \oplus \mathcal{O}_{\mathbb{P}^1}(-(b-e)),\mathcal{O}_{\mathbb{P}^1})$$ 

Hence $e$ is an isomorphism.

\color{red}
COMMENT: The following is just for our clarity:
We \textbf{probably} have a more general pairing,

$$\operatorname{Hom}((Rp_*T_Y(K))^*,\mathcal{O}_{\mathbb{P}^1}) \times \operatorname{Ext}^1((Rp_*T_Y(-H))^*,\mathcal{O}_{\mathbb{P}^1}) \to \operatorname{Ext}^2((Rp_*T_Y(-H+K))^*,\mathcal{O}_{\mathbb{P}^1}) $$

The above boils down to four pairings

$$\operatorname{Hom}((R^1p_*T_Y(K))^*,\mathcal{O}_{\mathbb{P}^1}) \times \operatorname{Ext}^1((p_*T_Y(-H))^*,\mathcal{O}_{\mathbb{P}^1}) \to \operatorname{Ext}^1((R^1p_*T_Y(-H+K))^*,\mathcal{O}_{\mathbb{P}^1}) $$

 which is the only non zero pairing. The second third and fourth pairings are between


$$\operatorname{Ext}^1((p_*T_Y(K))^*,\mathcal{O}_{\mathbb{P}^1}) \times \operatorname{Hom}((R^1p_*T_Y(-H))^*,\mathcal{O}_{\mathbb{P}^1}) \to \operatorname{Ext}^1((R^1p_*T_Y(-H+K))^*,\mathcal{O}_{\mathbb{P}^1}) $$

$$\operatorname{Hom}((R^1p_*T_Y(K))^*,\mathcal{O}_{\mathbb{P}^1}) \times \operatorname{Hom}((R^1p_*T_Y(-H))^*,\mathcal{O}_{\mathbb{P}^1}) \to \operatorname{Hom}((R^2p_*T_Y(-H+K))^*,\mathcal{O}_{\mathbb{P}^1}) $$

$$\operatorname{Ext}^1((p_*T_Y(K))^*,\mathcal{O}_{\mathbb{P}^1}) \times \operatorname{Ext}^1((p_*T_Y(-H))^*,\mathcal{O}_{\mathbb{P}^1}) \to \operatorname{Ext}^2((p_*T_Y(-H+K))^*,\mathcal{O}_{\mathbb{P}^1}) $$

and are all zero either because $R^1p_*T_Y(-H) = p_*T_Y(-H+K) = 0$ and there is no $R^2p_*$.




%Using the fact that both maps $b \circ a$ and $c$ naturally arise out the data of the ribbon $\widetilde{Y} \subset \mathbb{P}^M$, or in other words chasing through the sequences \ref{1}, \ref{2} and \ref{3}, it can be shown that $e = \operatorname{id}$.


\color{black}

Hence $a$ is surjective. This ends our proof.


\QEDB

\begin{theorem}\label{non-extendability of K3 surfaces}
    Let $X$ be a $K3$ surface with Picard group $\mathbb{Z}H$, where $H$ is the generator. Let $H^2 = 2g-2$. Let $X \subset \mathbb{P}^M$ be the embedding given by the complete linear series of $rH$. Assume either of the following conditions hold:
    \begin{enumerate}
        \item[(1)] $r \geq 5, g \geq 3$
        \item[(2)] $r = 4, g \geq 4$
        \item[(3)] $r = 3, g \geq 5$ 
        \item[(4)] $r = 2$, $g \geq 7$ 
        %and 
                  % \begin{enumerate}
                   %    \item[(a)]  $g = 2pq+1, p \geq 2, q \geq 2$, $p$, $q$ coprime or
                    %   \item[(b)] $g = 2pq-p^2+1, p \geq 2, q \geq p+1$, $p$, $q$ coprime 
                   %\end{enumerate}
        
        
        
        \item[(5)] $r = 1$, $g = 11$ or $g \geq 13$  
                   %\begin{enumerate}
                     %  \item[(a)] $g = 6q+1$, $3 \nmid q$, $q \geq 4$
                     %  \item[(b)] $g = 6q+1$, $3 \nmid q+1$, $q \geq 5$
                      % \item[(c)] $g = 8q+1$, $2 \nmid q$, $q \geq 4$
                      % \item[(d)] $g = 2pq+1$, $p, q$ coprime, $p \geq 5, q \geq 5$
                       %\item[(e)] $g = 2pq-p^2+1$, $p, q$ coprime, $p \geq 4, q \geq p+2$ 
                   %\end{enumerate}
        
    \end{enumerate}
\end{theorem}

Then $X \subset \mathbb{P}^M$ is not extendable.

\begin{proof}
    We first show that a general $K3$ surface of index $r$ and genus $g$ as mentioned in Theorem \ref{non-extendability of K3 surfaces} is not extendable. For this, we use Proposition \ref{calculating beta}, to show that $\beta = 0$ for $a \geq 3$ and $\gamma = 0$ for $a = 2$ and hence $\alpha(\widetilde{Y}) = 0$ for a suitably chosen $K3$ carpet $\widetilde{Y}$. 
    %type II (\textcolor{red}{you are referring to Kondoa's paper. Should we indicate what Type II mean...makes it less irrritating for the reader}) normal crossing $Y_1 \bigcup Y_2$. 
    \begin{enumerate}
        \item[(1)] For $r \geq 5$, we can choose $a = r, e = 0, b = rm, m \geq 1$ to get non-extendability of general $K3$ surfaces of index $r$ and genus $2m+1$, $m \geq 1$ and $a = r, e = 1, b = rm, m \geq 2$ to get non-extendability of general $K3$ surfaces of index $r$ and genus $2m$ with $m \geq 2$. 

        \smallskip
        
        \item[(2)] For $r = 4$, we can choose $a = 4, e = 0, b = 4m, m \geq 2$ to get non-extendability of general $K3$ surfaces of index $r$ and genus $2m+1$, $m \geq 2$ and $a = 4, e = 1, b = 4m, m \geq 2$ to get non-extendability of general $K3$ surfaces of index $4$ and genus $2m$ with $m \geq 2$. 

       \smallskip
        
        \item[(3)] For $r = 3$, we can choose $a = 3, e = 0, b = 3m, m \geq 2$ to get non-extendability of general $K3$ surfaces of index $3$ and genus $2m+1$, $m \geq 1$ and $a = 3, e = 1, b = 3m, m \geq 3$ to get non-extendability of general $K3$ surfaces of index $3$ and genus $2m$ with $m \geq 2$. 
        
        \smallskip

        \item[(4)] For $r = 2$, we use part $(2)$. One can choose $a = 2, e = 0, b = 2m, m \geq 3$ to get non-extendability of general $K3$ surfaces of index $2$ and genus $2m+1$, $m \geq 3$ and $a = 2, e = 1, b = 2m, m \geq 4$ to get non-extendability of general $K3$ surfaces of index $2$ and genus $2m$ with $m \geq 4$.
       % \item[(4)] For $r = 2$, we can choose $a = 2p, e = 0, b = 2q$ with $p, q$ coprime, $p \geq 2, q \geq 2$ to get non-extendability of general $K3$ surfaces of index $2$ and genus $2pq+1$ and $a = 2p, e = 1, b = 2q$ with $p, q$ coprime, $p \geq 2, q \geq p+1$ to get non-extendability of general $K3$ surfaces of index $2$ and genus $2pq-p^2+1$
        
       \smallskip

       \item[(5)]  For $r = 1$, we use part $(3)$. One can choose $a = 1, e = 0, b = m, m \geq 5$ to get non-extendability of general $K3$ surfaces of index $1$ and genus $2m+1$, $m \geq 5$ and $a = 1, e = 1, b = m, m \geq 7$ to get non-extendability of general $K3$ surfaces of index $1$ and genus $2m$ with $m \geq 7$.
        
      %  \item[(5)] For $r = 1$, we can choose $a = 3, e = 0, b = q$ with $3 \nmid q$ with $q \geq 4$ to get non-extendability of general $K3$ surfaces of index $1$ and genus $6q+1$ or $a = 3, e = 0, b = q$ with $3 \nmid q$ with $q \geq 6$ to get non-extendability of general $K3$ surfaces of index $1$ and genus $6q-7$ or $a = 4, e = 0, b = q$ with $2 \nmid q$ with $q \geq 4$ to get non-extendability of general $K3$ surfaces of index $1$ and genus $8q+1$ or $a = p, e = 0, b = q$ with $p, q$  with $p \geq 5, q \geq 3$, $p, q$ coprime to get non-extendability of general $K3$ surfaces of index $1$ and genus $2pq+1$ or $a = p, e = 1, b = q$ with $p, q$ coprime, $p \geq 4, q \geq p+2$ to get non-extendability of general $K3$ surfaces of index $1$ and genus $2pq-p^2+1$. 
        
    \end{enumerate}
\end{proof}


\section{Extendability of $K3$ surfaces of lower genus and classification of Fano threefolds of Picard rank $1$}

\subsection{Extendability of $K3$ surfaces of lower genus}
We give an upper bound to $\alpha(X)$ for a $K3$ surface $X \subset \mathbb{P}^N$ of index $r$ and genus $g$. Once again as in Theorem \ref{non-extendability of K3 surfaces}, using Theorem \ref{ZL for K3 carpets}, we compute $\alpha(\widetilde{Y})$ for a suitably chosen $K3$ carpet $\widetilde{Y}$, which by semicontinuity gives us an upper bound. We first state our main Theorem of this section

\begin{theorem}\label{ZL for lower genus $K3$ surfaces}

Let $X$ be a $K3$ surface with Picard group $\mathbb{Z}H$, where $H$ is the generator. Let $H^2 = 2g-2$. Let $X \subset \mathbb{P}^M$ be the embedding given by the complete linear series of $rH$. 
    \begin{enumerate}
        \item If $r = 1$, $g = 12$, then $\alpha(X) \leq 1$
        \item If $r = 1$, $g = 10$, then $\alpha(X) \leq 3$  
        \item If $r = 1$, $g = 9$, then $\alpha(X) \leq 4$
        \item If $r = 1$, $g = 8$, then $\alpha(X) \leq 6$
        \item If $r = 1$, $g = 7$, then $\alpha(X) \leq 8$
        \item If $r = 1$, $g = 6$, then $\alpha(X) \leq 10$ (\textcolor{blue}{exceptional case})
        \item If $r = 2$, $g = 3$, then $\alpha(X) \leq 10$ (\textcolor{blue}{exceptional case})
        %(\textcolor{blue}{Needs improvement by $1$.})
        \item If $r = 2$, $g = 4$, then $\alpha(X) \leq 6$ 
        \item If $r = 2$, $g = 5$, then $\alpha(X) \leq 3$ 
        \item If $r = 2$, $g = 6$, then $\alpha(X) \leq 1$ 
        \item If $r = 3$, $g = 3$, then $\alpha(X) \leq 4$
        \item If $r = 3$, $g = 4$, then $\alpha(X) \leq 1$ 
        \item If $r = 4$, $g = 3$, then $\alpha(X) \leq 1$
        
       
        
    \end{enumerate}
    
\end{theorem}

\begin{proof}
  We first prove for the case $r = 1$, by showing that in this case $\alpha({\widetilde{Y}}) \leq h^0(-H-2K_Y)$. We have $6 \leq g \leq 10$ or $g = 12$. We can degenerate such a $K3$ surface to a $K3$ carpet $\widetilde{Y} \subset \mathbb{P}^M$, extending $Y \subset \mathbb{P}^M$, where $Y = \mathbb{F}_e$ embedded by the complete linear series of $C_0+bf$ where $e = 0$ or $e = 1$ depending on whether $g$ is odd or even respectively. Such a $\widetilde{Y}$ has a canonical ribbon section with arithmetic genus $p_a = g = 2b-e+1$. So it is enough to find an upper bound for $\alpha(\widetilde{Y})$ for a $K3$ carpet $\widetilde{Y} \subset \mathbb{P}^M$, extending $Y \subset \mathbb{P}^M$, where $Y = \mathbb{F}_e$ embedded by the complete linear series of $C_0+bf$ where the pair $(e,b)$ is one of $(0,3), (0,4), (1,3), (1,4), (1,5), (1,6)$. To see this, note that if $b \geq e+3$ by equations \ref{42}-\ref{62}, we have (since the only vanishing, we do not have is that of $h^0(-H-2K_Y)$ for which we require $b \geq 2e+5$)  the following modified version of equation \ref{63}   
  
 % \ref{50}, we have $H^0(N_{\widetilde{Y}/\mathbb{P}^M}(-H+K_Y)) = 0, H^2(N_{\widetilde{Y}/\mathbb{P}^M}(-H+K_Y)) = 0$ for $b \geq e+3$. Hence for $b \geq e+3$, we have the following weaker version of equation \ref{63}
\begin{equation}\label{70}
    h^0(N_{\widetilde{Y}/\mathbb{P}^M}(-\widetilde{H})) = N+1 + h^1(N_{\widetilde{Y}/\mathbb{P}^M}(-\widetilde{H}))-h^1(-H-2K_Y) + h^0(-H-2K_Y)
\end{equation}

And then we have shown in the proof of part $(3)$ of Theorem \ref{calculating beta} (right after equation \ref{63}) that for $b \geq e+3$, $h^1(N_{\widetilde{Y}/\mathbb{P}^M}(-\widetilde{H}))-h^1(-H-2K_Y)$. Hence by equation \ref{70}, we have that for $b \geq e+3$, $\alpha({\widetilde{Y}}) \leq h^0(-H-2K_Y)$. Now computing $h^0(-H-2K_Y)$ for the cases $(0,3), (0,4), (1,3), (1,4), (1,5), (1,6)$, give us the bounds as stated in the Theorem.  \par

We now prove for $r = 2$. Since $3 \leq g \leq 6$. We can degenerate such a $K3$ surface to a $K3$ carpet $\widetilde{Y} \subset \mathbb{P}^M$, extending $Y \subset \mathbb{P}^M$, where $Y = \mathbb{F}_e$ embedded by the complete linear series of $2C_0+bf$, where $b = 2m$ and $e = 0$ or $e = 1$ depending on whether $g$ is odd or even respectively. In this case $g = 2m-e+1$. So it is enough to find an upper bound for $\alpha(\widetilde{Y})$ for a $K3$ carpet $\widetilde{Y} \subset \mathbb{P}^M$, extending $Y \subset \mathbb{P}^M$, where $Y = \mathbb{F}_e$ embedded by the complete linear series of $2C_0+2mf$ where the pair $(e,m)$ is one of $(0,1), (0,2), (1,2), (1,3)$. We have shown in Theorem \ref{ZL for K3 carpets}, part $(2)$, that $$\alpha(\widetilde{Y})  \leq  \gamma \leq h^1(T_Y(-H+K_Y))+h^0(N_{Y/\mathbb{P}^M}(-H)+h^0(-H-2K_Y)-M-1 $$
So we calculate the right hand side of the above equation for the above mentioned values of $(e,m)$. By Lemma \ref{T_Y(-H)}, $(2)$, we have that for $m \geq e+1/2$, $h^1(T_Y(-H+K_Y)) = 0$. Now in the proof of part $(2)$ of Theorem \ref{calculating beta}, using equations \ref{35}-\ref{41}, we have shown that if $m \geq e/2+3/2$, $h^0(N_{Y/\mathbb{P}^M}(-H) = M+1$. Hence except for the case for the cases $(e,m) = (0,1)$, $\alpha(\widetilde{Y}) \leq h^0(-H-2K_Y)$. This gives us the upper bounds as stated in the statement for the cases $r = 2, 4 \leq g \leq 6$. On the other hand once again in the proof of part $(2)$ of Theorem \ref{calculating beta}, in equation \ref{38}, we have that $H^1(p^*(T_{\mathbb{P}^1})(-H)) = H^1(\mathcal{O}_{\mathbb{P}^1}(-2) = \mathbb{C}$, which in turn gives us $h^0(N_{Y/\mathbb{P}^M}(-H) = M+2$ from the sequence in \ref{35}. Hence in this case $\alpha(\widetilde{Y}) \leq h^0(-H-2K_Y)+1$, which gives us the remaining case $r = 1, g = 3$. \par

We now prove for $r = 3,4$. We can degenerate such a $K3$ surface to a $K3$ carpet $\widetilde{Y} \subset \mathbb{P}^M$, extending $Y \subset \mathbb{P}^M$, where $Y = \mathbb{F}_e$ embedded by the complete linear series of $rC_0+rmf$, where $e = 0$ or $e = 1$ depending on whether $g$ is odd or even respectively. Again we have $g = 2m-e+1$. So it is enough to find an upper bound for $\alpha(\widetilde{Y})$ for a $K3$ carpet $\widetilde{Y} \subset \mathbb{P}^M$, extending $Y \subset \mathbb{P}^M$, where $Y = \mathbb{F}_e$ embedded by the complete linear series of $rC_0+rmf$ where the tuple $(r,e,m)$ is one of $(3,0,1), (3,1,2), (4,0,1)$. In each case, we calculate $\beta$ in Theorem \ref{ZL for K3 carpets} which gives an upper bound for $\alpha(\widetilde{Y})$. In each case, one can check that the only contribution to $\beta$ comes from $h^0(-H-2K_Y)$, which gives us the stated bounds.
\end{proof}

\subsection{Classification of Fano-threefolds}
\begin{proposition}\label{H^0(N(-2))}
Let $\widetilde{Y} \hookrightarrow \mathbb{P}^M$ be a $K3$ carpet embedded by the complete linear series of a very ample line bundle $\widetilde{H}$ such that the embedding induced on the reduced part is $Y = \mathbb{F}_e \hookrightarrow \mathbb{P}^N$ embedded by the complete linear series of a very ample line bundle $H = aC_0+bf$. Assume $H^1(-H+K_Y) = 0$. Then $$h^0(N_{\widetilde{Y}/\mathbb{P}^N}(-2\widetilde{H})) \leq h^1(T_Y(-2H+K_Y)) + h^1(T_Y(-2H)) - h^0(T_Y(-2H)) +  h^1(-2H-K_Y) + h^0(-2H-2K_Y) $$
Consequently, $h^0(N_{\widetilde{Y}/\mathbb{P}^N}(-2\widetilde{H})) = 0$ if either of the following conditions hold
\begin{enumerate}
    \item $a = 2, b \geq 3, e = 0$
    \item $a \geq 3, b \geq 2, e = 0$
    \item $a \geq 2, b \geq a/2+1, e = 1$
    \item $a \geq 2, b \geq ae/2+1/2, e \geq 2$
    
\end{enumerate}
Moreover, if $a = 2, b = 2, e = 0$, $h^0(N_{\widetilde{Y}/\mathbb{P}^N}(-2\widetilde{H})) \leq 1$
\end{proposition}

\begin{proof}
    We follow the proof of \cite[Theorem $2.2$]{BM24}, and indicate the changes. We take $L = K_Y$, tensor by $-2\widetilde{H}$ instead of $-\widetilde{H}$ and follow the proof verbatim. Equations $2.1-2.9$ in \cite[Theorem $2.2$]{BM24}, yields 
    $$h^0(N_{\widetilde{Y}/\mathbb{P}^M}(-2\widetilde{H})) \leq  h^0(N_{Y/\mathbb{P}^M}(-2H+K_Y)) + h^0(N_{Y/\mathbb{P}^M}(-2H)) +  h^1(-2H-K_Y) + h^0(-2H-2K_Y)$$
Now tensoring the exact sequence $2.10-2.11$ in \cite[Theorem $2.2$]{BM24} by $-2H$ now yields 
$$h^0(N_{Y/\mathbb{P}^M}(-2H)) \leq h^1(T_Y(-2H)) - h^0(T_Y(-2H))$$
Similarly tensoring $2.10-2.11$ in \cite[Theorem $2.2$]{BM24} by $-2H+K_Y$ yields 
$$h^0(N_{Y/\mathbb{P}^M}(-2H+K_Y)) \leq h^1(T_Y(-2H+K_Y))$$
Hence we have 
$$h^0(N_{\widetilde{Y}/\mathbb{P}^N}(-2\widetilde{H})) \leq h^1(T_Y(-2H+K_Y)) + h^1(T_Y(-2H)) - h^0(T_Y(-2H)) +  h^1(-2H-K_Y) + h^0(-2H-2K_Y) $$
Now the rest of the proof follows from \Cref{-H-K_Y}, \Cref{-H-2K_Y}, \Cref{T_Y(-H)}. For the ease of the reader we point out that in the case $a = 2, b = 2, e = 0$, the only term that contributes in the upper bound is $h^0(-2H-2K_Y) = 1$. 
\end{proof}


\begin{proposition}\label{r = 1, H^0(N(-2))}
Let $\widetilde{Y} \hookrightarrow \mathbb{P}^M$ be a $K3$ carpet embedded by the complete linear series of a very ample line bundle $\widetilde{H}$ such that the embedding induced on the reduced part is $Y = \mathbb{F}_e \hookrightarrow \mathbb{P}^N$ embedded by the complete linear series of a very ample line bundle $H = aC_0+bf$. Assume $H^1(-H+K_Y) = 0$. Then $$h^0(N_{\widetilde{Y}/\mathbb{P}^N}(-2\widetilde{H})) \leq h^1(T_Y(-2H+K_Y)) + h^0(N_{Y/\mathbb{P}^M}(-2H))+ h^0(-2H-2K_Y) $$
Consequently, $h^0(N_{\widetilde{Y}/\mathbb{P}^N}(-2\widetilde{H})) = 0$ if $a = 1, b \geq e+5/2$.
Moreover, if $a = 1, b = 3, e = 1$, $h^0(N_{\widetilde{Y}/\mathbb{P}^N}(-2\widetilde{H})) \leq 1$
\end{proposition}


\begin{definition}\label{Hilbert scheme of Fanos and K3}
    Let $\mathcal{V}_{r,g,1}$ denote the Hilbert scheme of smooth Fano threefolds of Picard rank one, index $r$ and genus $g$. Similarly let $\mathcal{H}_{r,g}$ denote the Hilbert scheme of $K3$ surfaces of index $r$ and genus $g$, i.e, $K3$ surfaces $X \subset \mathbb{P}^{1+r^2(g-1)}$ embedded the complete linear series of a very ample line bundle $rH$, where $H^2 = 2g-2$.
\end{definition}


\begin{corollary}\label{upper bound on cone over $K3$ surface}

Let $X$ be a $K3$ surface with Picard group $\mathbb{Z}H$, where $H$ is the generator. Let $H^2 = 2g-2$. Let $X \subset \mathbb{P}^M$ be the embedding given by the complete linear series of $rH$. 

\begin{enumerate}
    \item If $r = 1, g \geq 7$ or $r \geq 2, g \geq 4$. Then the dimension of the tangent space of $\mathcal{V}_{r,g,1}$ at the point $C(X) \subset \mathbb{P}^{M+1}$ is given by 

\begin{equation}\label{upper bounds}
    \operatorname{dim}T_{C(X)}(\mathcal{V}_{r,g,1}) = h^0(N_{X/\mathbb{P}^M})+h^0(N_{X/\mathbb{P}^{M}}(-1)) = 18+(2+r^2(g-1))^2+\alpha(X)+(2+r^2(g-1))
\end{equation}

   \item If $r = 1, g = 6$ or $r = 2, g = 3$. Then the dimension of the tangent space of $\mathcal{V}_{r,g,1}$ at the point $C(X) \subset \mathbb{P}^{M+1}$ is given by 

\begin{equation}\label{upper bounds for r = 2, g = 3}
    \operatorname{dim}T_{C(X)}(\mathcal{V}_{r,g,1}) = h^0(N_{X/\mathbb{P}^M})+h^0(N_{X/\mathbb{P}^{M}}(-1)) + 1 = 18+(2+r^2(g-1))^2+\alpha(X)+(2+r^2(g-1))+1
\end{equation}


\end{enumerate}

    
\end{corollary}

\begin{lemma}\label{hyperplane section of a general fano is a general K3}

Let $V$ be a general element of $\mathcal{V}_{r,g,1}$. Then a general hyperplane section of $V$ corresonds to a general $K3$ surface in $\mathcal{H}_{r,g}$.
    
\end{lemma}

\begin{proof}
    The proof follows from standard deformation theory arguments (see for example \cite[Lemma $3.7$]{CLM98}).
\end{proof}

\begin{theorem}\label{classification of fanos with index two}
Let $\mathcal{V}_{r,g,1}$ denote the Hilbert scheme of smooth Fano threefolds of Picard rank one, index $r \geq 2$ and genus $g \geq 3$ or $r = 1, g \geq 6$. Then 
\begin{enumerate}
    \item $\mathcal{V}_{r,g,1} = \varnothing$ if $(r = 1, g = 11)$, $(r = 1, g \geq 13)$, $(r = 2, g \geq 7)$, $(r = 3, g = 3)$, $(r = 3, g \geq 5)$, $(r = 4, g \geq 4)$ or $(r = 5, g \geq 3)$. 
    \item For $(r = 1, 6 \leq g \leq 10)$, $(r = 1, g = 12)$, $(r = 2, 3 \leq g \leq 6), (r = 3, g = 4), (r = 4, g = 3)$, $\mathcal{V}_{r,g,1}$ is irreducible, and the families of Fano and Iskovskih form an open dense irreducible subset of $\mathcal{V}_{r,g,1}$.
\end{enumerate}
\end{theorem}

\begin{proof}
  \noindent\textit{Proof of $(1)$} First note that for every $r$, there is a unique component $\mathcal{H}_{r,g}$ of the Hilbert scheme parameterizing $K3$ surfaces $X \subset \mathbb{P}^{h^0(rH)-1}$ embedded by the complete linear of a very ample line bundle $rH$, where $H^2 = 2g-2$. Now excepting the case, $(r = 3, g = 3)$, part $(1)$ follows from \Cref{non-extendability of K3 surfaces}. For the case $(r=3, g=3)$, note that if the hyperplane section of a Fano threefold $V$ in $\mathcal{V}_{r,g,1}$ be denoted by $rL$ and $L^3 = d$, then computing $h^0(rL) = h^0(rH)-1$ by Riemann-Roch theorem, we have $d = 2(g-1)/r$. But for $r = 3, g = 3$, $d$ cannot be an integer. This resolves the case. \par
  \textit{Proof of $(2)$} First one notes that for each $r$ and $g$, $\mathcal{H}_{r,g}$ is irreducible. Now the list of examples of Fano-Iskovskih-Mori-Mukai (see for example \cite[Table $3.1$]{CLM98}), show that in each of the cases mentioned by part $(2)$, there exist $K3$ surfaces in $\mathcal{H}_{r,g}$ which are extendable to smooth Fano threefolds. Hence by \Cref{hyperplane section of a general fano is a general K3}, in each of the cases, $\mathcal{H}_{r,g}$ is dominated by the unique irreducible component $\mathcal{V}_{r,g,1}$ containing the examples of Fano-Iskovskih-Mori-Mukai. We end the proof by showing, this is the unique irreducible component of $\mathcal{V}_{r,g,1}$ and that the known examples form an open dense subset of $\mathcal{V}_{r,g,1}$.  To see this, we observe first as in \cite{CLM93} and \cite{CLM98}, that a projectively Cohen-Macaulay scheme degenerates along a flat family to the cone over its hyperplane section. Hence once again using, \Cref{hyperplane section of a general fano is a general K3}, it is enough to show that for a general $K3$ surface $X \subset \mathbb{P}^{M}$ in $\mathcal{H}_{r,g}$, the cone $C(X) \subset \mathbb{P}^{M+1}$ is a smooth point of $\mathcal{V}_{r,g,1}$. For the stated values of pairs $(r,g)$, 
 % \Cref{upper bounds} and \Cref{upper bounds for r = 2, g = 3} in 
  \Cref{upper bound on cone over $K3$ surface} along with \Cref{ZL for lower genus $K3$ surfaces}, gives an upper bound to $\operatorname{dim}T_{C(X)}(\mathcal{V}_{r,g,1})$ which is listed as follows 
  \begin{enumerate}
      \item $\operatorname{dim}T_{C(X)}(\mathcal{V}_{1,6,1}) \leq 85$

      \item $\operatorname{dim}T_{C(X)}(\mathcal{V}_{1,7,1}) \leq 98$

      \item $\operatorname{dim}T_{C(X)}(\mathcal{V}_{1,8,1}) \leq 114$

      \item $\operatorname{dim}T_{C(X)}(\mathcal{V}_{1,9,1}) \leq 132$
      
      \item $\operatorname{dim}T_{C(X)}(\mathcal{V}_{1,10,1}) \leq 153$

      \item $\operatorname{dim}T_{C(X)}(\mathcal{V}_{1,12,1}) \leq 201$
      
      \item $\operatorname{dim}T_{C(X)}(\mathcal{V}_{2,3,1}) \leq 139$

      \item $\operatorname{dim}T_{C(X)}(\mathcal{V}_{2,4,1}) \leq 234$

      \item $\operatorname{dim}T_{C(X)}(\mathcal{V}_{2,5,1}) \leq 363$

      \item $\operatorname{dim}T_{C(X)}(\mathcal{V}_{2,6,1}) \leq 525$

      \item $\operatorname{dim}T_{C(X)}(\mathcal{V}_{3,4,1}) \leq 889$

      \item $\operatorname{dim}T_{C(X)}(\mathcal{V}_{4,3,1}) \leq 1209$
  \end{enumerate}
Now for a smooth Fano threefold $V \in \mathcal{V}_{r,g,1}$, $V$ represents a smooth point of $\mathcal{V}_{r,g,1}$, since embedded deformations of an anticanonically embedded Fano variety of dimension at least $3$ can be shown to be unobstructed using the vanishing of $h^1(N_{V/\mathbb{P}^{M+1}})$ by Kodaira-Nakano vanishing theorem. Therefore the dimension of $\mathcal{V}_{r,g,1}$ at a smooth point is given by $\operatorname{dim}T_V(\mathcal{V}_{r,g,1}) = h^0(N_{V/\mathbb{P}^{M+1}})$. Now using the descriptions of the examples of Fano-Iskovskih-Mori-Mukai, one can calculate the dimension $h^0(N_{V/\mathbb{P}^{M+1}})$ (see for example \cite[Table $2$, pg $666$]{CLM93} and \cite[Table $3.1$]{CLM98}) to check that in each case, $\operatorname{dim}T_V(\mathcal{V}_{r,g,1})$ is exactly the upper bound just obtained. Hence $C(X)$ is also a smooth point of $\mathcal{V}_{r,g,1}$. 

  
\end{proof}


\vspace{0.5cm}

\bibliographystyle{plain}
\begin{thebibliography}{KMM92}
\color{black}

\bibitem[AFS13]{AFS13}
Alper, Jarod; Fedorchuk, Maksym; Smyth, David Ishii
\emph{Finite Hilbert stability of (bi)canonical curves.}
Invent. Math.191(2013), no.3, 671–718.

\bibitem[Ale15]{Alexeev}
Alexeev, Valery 
\emph{Moduli of weighted hyperplane arrangements.}
Edited by Gilberto Bini, Martí Lahoz, Emanuele Macrì and Paolo Stellari, 
Adv. Courses Math. CRM Barcelona
Birkhäuser/Springer, Basel, 2015. 

\bibitem[AHL10]{AHL10}
Artebani, Michela; Hausen, Jürgen; Laface, Antonio
\emph{On Cox rings of K3 surfaces}
Compos. Math. 146 (2010), no. 4, 964–998.

\bibitem[BE95]{BE95}
Bayer, Dave; Eisenbud, David
\emph{Ribbons and their canonical embeddings}
Trans. Amer. Math. Soc. 347 (1995), no. 3, 719–756.

\bibitem[Bel24]{Bel24}
Belmans, Peter
\emph{Fanography: a tool to visually study the geography of Fano 3-folds.}
https://www.fanography.info

\bibitem[BGG23]{BGG} Bangere, Purnaprajna; Gallego, Francisco Javier; González, Miguel \emph{Deformations of hyperelliptic and generalized hyperelliptic polarized varieties.}
Mediterr. J. Math. 20 (2023), no. 2, Paper No. 83.
\color{black}

\bibitem[BM24]{BM24} 
Bangere, Purnaprajna; Mukherjee, Jayan  
\emph{Extendability of projective varieties via degeneration to ribbons with applications to Calabi-Yau threefolds} (arXiv:2409.03960)


\bibitem[BGM24]{BGM24} Bangere, Purnaprajna; Gallego, Francisco Javier; Mukherjee; Jayan \emph{Projective degenerations of $K3$ surfaces to a union of Hirzebruch surfaces not embedded as scrolls.} (in preparation)

\bibitem[BGM24b]{BGM24b} Bangere, Purnaprajna; Gallego, Francisco Javier; Mukherjee; Jayan \emph{Projective smoothing of simple normal crossing projective bundles.} (in preparation)

\bibitem[BGM24c]{BGM24c} Bangere, Purnaprajna; Gallego, Francisco Javier; Mukherjee; Jayan \emph{Projective smoothing of simple normal crossing complete intersections.} (in preparation)


\bibitem[Bea02]{B02}
Beauville, A. \emph{Fano threefolds and $K3$ surfaces}
https://arxiv.org/abs/math/0211313

\bibitem[CLM93]{CLM93}
Ciliberto, C; Lopez, A; Miranda, R
\emph{Projective degenerations of $K3$ surfaces, Gaussian maps and Fano threefolds}
Invent. Math. 114, 641 667 (1993)

\bibitem[CLM98]{CLM98}
Ciliberto, Ciro; Lopez, Angelo Felice; Miranda, Rick
\emph{Classification of varieties with canonical curve section via Gaussian maps on canonical curves.}
Amer. J. Math.   120 (1998), no. 1, 1–21.

\bibitem[CLS10]{CLS}
Cox, David A.; Little, John B.; Schenck, Henry K.
\emph{Toric varieties}
Grad. Stud. Math., 124
American Mathematical Society, Providence, RI, 2011. xxiv+841 pp
\url{https://www.mimuw.edu.pl/~jarekw/pragmatic2010/CoxLittleSchenckJan2010.pdf}
\url{https://www.researchgate.net/profile/Hal-Schenck/publication/228552810_Toric_Varieties_Chapters_8-11/links/0046351f1e7da63620000000/Toric-Varieties-Chapters-8-11.pdf}

\color{black}
\bibitem[Deb01]{Debarre} Debarre, O., \emph{Higher-dimensional algebraic geometry}, 
Universitext Springer--Verlag, New York, 2001.
\color{black}

\bibitem[Deo18]{Deo18}
Deopurkar, Anand
\emph{The canonical syzygy conjecture for ribbons}
Math. Z. 288 (2018), no. 3-4, 1157–1164.

\bibitem[DFS16]{DFS16}
Deopurkar, Anand; Fedorchuk, Maksym; Swinarski, David
\emph{Toward GIT stability of syzygies of canonical curves}
Algebr. Geom. 3 (2016), no. 1, 1–22.


\bibitem[DM69]{DM69}
Deligne, P.; Mumford, D.
\emph{The irreducibility of the space of curves of given genus.}
Inst. Hautes Études Sci. Publ. Math.(1969), no.36, 75–109.

\bibitem[EH00]{EH} Eisenbud, David; Harris, Joe 
\emph{The geometry of schemes.}
Grad. Texts in Math., 197
Springer-Verlag, New York, 2000.

\bibitem[Fed18]{Fed18}
Fedorchuk, Maksym
\emph{Geometric invariant theory of syzygies, with applications to moduli spaces.}
Geometry of moduli, 107–134. Abel Symp., 14 Springer, Cham, 2018

\bibitem[FJ13]{FJ13}
Fedorchuk, Maksym; Jensen, David
\emph{Stability of 2nd Hilbert points of canonical curves.}
Int. Math. Res. Not. IMRN(2013), no.22, 5270–5287.

\bibitem[Fri83]{F83}
Friedman, Robert; 
\emph{Global smoothings of varieties with normal crossings.} 
Ann. of Math. (2) 118 (1983), no. 1, 75--114.

 

\color{black}

\bibitem[Fuj14a]{KFujita} Fujita, Kento, \emph{Simple normal crossing Fano varieties and log Fano manifolds.}
Nagoya Math. J. 214 (2014), 95-123.

\bibitem[Fuj14b]{Fujino} Fujino, Osamu.
\emph{Fundamental theorems for semi log canonical pairs.} Algebr. Geom. 1 (2014), no.2, 194-228.

\color{black}

\bibitem[GGP08]{GGP08}
Gallego, Francisco Javier; González, Miguel; Purnaprajna, Bangere P., 
\emph{$K3$ double structures on Enriques surfaces and their smoothings.}
J. Pure Appl. Algebra 212 (2008), no. 5, 981--993
\bibitem[GGP13]{GGP13}
Gallego, F.J., González, M., Purnaprajna, B.P., 
\emph{An infinitesimal condition to smooth ropes.}
Rev Mat Complut 26, 253–269 (2013).

\bibitem[GGP14]{GGP14}
Gallego, F.J.; González, Miguel; Purnaprajna, B.P., 
\emph{Smoothable locally non Cohen–Macaulay multiple structures on curves.} 
Collect. Math. 65, 417–433 (2014).

\bibitem[GP97]{GP97}
Gallego, F.J., Purnaprajna, B.P.
\emph{Degenerations of $K3$ surfaces in projective space}
Transactions of the American Mathematical Society
Volume 349, Number 6, June 1997, Pages 2477–2492

\bibitem[Gon06]{Go}
  Gonz\'alez, Miguel.
  \emph{Smoothing of ribbons over curves}.
  J. Reine Angew. Math. 591 (2006), 201--235.

  \color{black}

  \bibitem[Gro61]{EGA} A. Grothendieck, \emph{\'El\'ements 
de g\'eom\'etrie 
alg\'ebrique. III. \'Etude cohomologique des faisceaux 
coh\'erents. I.}, Inst. 
Hautes \'Etudes Sci. Publ. Math.  
{11}, (1961).

\bibitem[Har77]{Hartshorne}  
Hartshorne, R; \emph{Algebraic geometry}, 
Grad. Texts in Math., No. 52
Springer-Verlag, New York-Heidelberg, 1977

\color{black}
  
\bibitem[Har10]{H10}
Hartshorne, R;
\emph{Deformation Theory}
Springer-Verlag New York 2010, https://doi.org/10.1007/978-1-4419-1596-2

\bibitem[HH09]{HH09}
Hassett, Brendan; Hyeon, Donghoon
\emph{Log canonical models for the moduli space of curves: the first divisorial contraction.}
Trans. Amer. Math. Soc.361(2009), no.8, 4471–4489.

\bibitem[HH13]{HH13}
Hassett, Brendan; Hyeon, Donghoon
\emph{Log minimal model program for the moduli space of stable curves: the first flip.}
Ann. of Math. (2)177(2013), no.3, 911–968.

\bibitem[HV85]{HV} K. Hulek \& A. Van de Ven, 
\emph{The Horrocks-Mumford bundle and the Ferrand construction.}
Manuscripta Math. 50 (1985), 313-335.

\bibitem[Isk77]{I77}
Iskovskih, V. A.
\emph{Fano threefolds. I.}
(Russian)Izv. Akad. Nauk SSSR Ser. Mat.41 (1977), no.3, 516–562, 717.

\bibitem[Isk78]{I78}
Iskovskih, V. A.
\emph{Fano threefolds. II.}
(Russian)Izv. Akad. Nauk SSSR Ser. Mat.42 (1978), no.3, 506–549.

\bibitem[KN94]{KN94}
Kawamata, Y., Namikawa, Y. 
\emph{Logarithmic deformations of normal crossing varieties and smoothing of degenerate Calabi-Yau varieties.}
Invent Math 118, 395–409 (1994).

\bibitem[Kim17]{Ki17}
Kimura, Y.;
\emph{Structure of stable degeneration of $K3$ surfaces into pairs of rational elliptic surfaces}
JHEP03(2018)045, https://doi.org/10.48550/arXiv.1710.04984

\bibitem[Kle81]{KL81}
Kleppe, J.; 
\emph{The Hilbert Flag Scheme, its properties, and its connection with the Hilbert Scheme. Applications to curves in three-space. Thesis. University of Oslo, Preprint No. 5 (1981)}
\bibitem[Kle87]{KL87}
Kleppe, J.;
\emph{Non-reduced components of the Hilbert scheme of smooth space curves. In: Space Curve.}
Rocca di papa, 1985. Ghione, F., Peskine, C., Sernesi, E. (eds.) Springer Lect. Notes Math. No. 1266 (1987), 181-207

\bibitem[Kol10]{Kol10}
Koll\'ar. J.;
\emph{Moduli of varieties of general type.} arXiv:1008.0621(2010)

  \color{black}

\bibitem[Kol13]{Kol13}
Koll\'ar. J.; \emph{Singularities of the Minimal Model Program}, Cambridge Tracts in Math., 200
Cambridge University Press, Cambridge, 2013

  \color{black}

\bibitem[Kol14]{Kol14}
Koll\'ar, J.;
\emph{Simple normal crossing varieties with prescribed dual complex.}
Algebr. Geom. 1(2014), no.1, 57–68.

 \color{black}

\bibitem[KMM92]{KMM} Kollár, János; Miyaoka, Y.; Mori, S., 
\emph{Rational connectedness and boundedness of Fano manifolds.} J. Differential Geom. 36 (1992), no.3, 765--779.

 \color{black}

\bibitem[KSB88]{KSBA88}
Koll\'ar, J;  Shepherd-Barron, N.I.;
\emph{Threefolds and deformations of surface
singularities}
Invent. Math. 91 (1988), no. 2, 299–338.

\bibitem[KX20]{KX20}
Kollár, János; Xu, Chen Yang
\emph{Moduli of polarized Calabi-Yau pairs.}
Acta Math. Sin. (Engl. Ser.)36(2020), no.6, 631–637.



\bibitem[Kon85]{Ko85}
Kondo, S
\emph{Type II degenerations of $K3$ surfaces}
Nagoya Math. J. Vol. 99 (1985), 11-30



\bibitem[Kul77]{Ku77}
Kulikov, V. 
\emph{Degenerations of K3 surfaces and Enriques surfaces} Math. USSR Izvestiya 11 (1977), 957-989.





\color{black}

\bibitem[Laz04]{positivity}
Lazarsfeld, Robert. 
\emph{Positivity in algebraic geometry. I},
Ergeb. Math. Grenzgeb. (3), 48, 
Springer-Verlag, Berlin, 2004.

\bibitem[LR12]{LR}
Liu, Wenfei; Rollenske, Sönke. \emph{Two-dimensional semi-log-canonical hypersurfaces}, 
Matematiche (Catania) 67 (2012), no. 2, 185–202.

\color{black}
\bibitem[Man03]{Man03}
Manolache, N
\emph{Cohen Macaulay Nilpotent Schemes}
https://arxiv.org/abs/math/0312514

\bibitem[MM82]{MM82}
Mori, Shigefumi; Mukai, Shigeru
\emph{Classification of Fano 3-folds with $B2 \geq 2$.}
Manuscripta Math.36(1981/82/1982), no.2, 147–162.


\bibitem[Muk88]{M88}
Mukai, S
\emph{Curves, K3 Surfaces and Fano 3-folds of Genus $\leq$ 10}
Algebraic Geometry and Commutative Algebra, Academic Press,
1988, Pages 357-377.

  \color{black}
  
\bibitem[MFK94]{MFK} D. Mumford, J. Fogarty, F. Kirwan, 
\emph{Geometric Invariant Theory, third enlarged edition}, 
Springer-Verlag, 1994, Ergebnisse der Mathematik und ihrer Grenzgebiete, 
{34}.

  \color{black}
\bibitem[Par91]{Par91}
Pardini, Rita
\emph{Abelian covers of algebraic varieties.}
J. Reine Angew. Math.417(1991), 191–213.

\bibitem[PMR21]{PMR21}
Bangere, Purnaprajna; Mukherjee, Jayan; Raychaudhury, Debaditya
\emph{K3 carpets on minimal rational surfaces and their smoothings.}
Internat. J. Math.32(2021), no.6, Paper No. 2150032, 20 pp.


\bibitem[PP81]{PP81}
Persson, U; Pinkham, H 
\emph{Degenerations of surfaces with trivial canonical bundle}
Ann. of Math. 113 (1981), 45-66

\bibitem[RS22]{RS22}
Raicu, Claudiu; Sam, Steven V.
\emph{Bi-graded Koszul modules, K3 carpets, and Green's conjecture.}
Compos. Math.158(2022), no.1, 33–56.


\bibitem[Ser06]{Ser}
Sernesi, Edoardo.
\emph{Deformations of Algebraic Schemes.}
Grundlehren der Mathematischen Wissenschaften [Fundamental Principles of Mathematical Sciences], 334. Springer-Verlag, Berlin, 2006.

\bibitem[Sei92]{Sei92}
Seiler, Wolfgang K.
\emph{Deformations of ruled surfaces.}
J. Reine Angew. Math.426(1992), 203–219.

\bibitem[SP]{SP}
Stacks project authors
\emph{The Stacks project}
\url{https://stacks.math.columbia.edu}, 2023

\bibitem[Tho13]{T13}
Thompson, A;
\emph{Degenerations of $K3$ surfaces of degree two}
Transactions of the American Mathematical Society
Volume 366, Number 1, January 2014, Pages 219–243
S 0002-9947(2013)05759-5

\bibitem[Tot23]{Tot23}
Totaro, Burt;
\emph{Bott Vanishing for Fano threefolds}
\url{https://doi.org/10.48550/arXiv.2302.08142}


\bibitem[Tzi10]{TN10}
Tziolas, Nikolaos; 
\emph{Smoothings of schemes with nonisolated singularities.}
Michigan Math. J. 59 (2010), no. 1, 25--84


\bibitem[Tzi15]{TN15}
Tziolas, N; (2015). 
\emph{Smoothings of Fano Varieties With Normal Crossing Singularities.} 
Proceedings of the Edinburgh Mathematical Society, 58(3), 787-806. doi:10.1017/S0013091515000024

\bibitem[Vie95]{Vie95}
Viehweg, Eckart;
\emph{Quasi-projective moduli for polarized manifolds.}
Ergeb. Math. Grenzgeb. (3), 30[Results in Mathematics and Related Areas (3)] Springer-Verlag, Berlin, 1995. viii+320 pp.
ISBN:3-540-59255-5



 






  

\end{thebibliography}

\end{document}

